% =====================================================================
%  Εξέλιξη και Σύγκριση Πρωτοκόλλων Συναίνεσης σε Κατανεμημένα Συστήματα
%  Παναγιώτης Δαμπάνης — 711070095
%  Τμήμα Μηχανικών Πληροφορικής και Υπολογιστών, ΠΑΔΑ
%
%  ΑΥΤΗ Η ΕΚΔΟΣΗ ΑΚΟΛΟΥΘΕΙ ΤΗΝ ΑΓΓΛΙΚΗ. Το κείμενο γράφεται στο
%  content/en/*.tex και μεταφράζεται στο content/el/*.tex. Το «make sync»
%  δείχνει ποια κομμάτια έχουν μείνει πίσω. Δες docs/SYNC.md.
%
%  >>> ΜΕΤΑΓΛΩΤΤΙΣΕ ΜΕ XeLaTeX <<<  (Overleaf: Menu -> Compiler -> XeLaTeX)
%
%  Η σειρά των σελίδων είναι του Προτύπου ΤΜΠΥ (docs/PROTYPO-TMPY.docx).
%  Αν την αλλάξεις, άλλαξέ την και στο docx/assemble.py — ο έλεγχος V07
%  συγκρίνει τα δύο.
% =====================================================================
\documentclass[11pt,a4paper,oneside]{report}

% =====================================================================
%  common.tex — everything that is the same in both languages.
%
%  Compile with XeLaTeX. pdfLaTeX will not do: the Greek edition needs
%  fontspec, and both editions must use the same font files so the two
%  PDFs look like the same document.
%
%  Language-specific setup lives in style/lang-en.tex and lang-el.tex.
% =====================================================================

% ---- Fonts ----------------------------------------------------------
\usepackage{fontspec}
\usepackage{polyglossia}

% The department requires Times New Roman at 11 pt (Guidelines A.2).
% Times New Roman itself is not redistributable, so the build asks for it
% first and falls back to Liberation Serif, which is metric-compatible:
% same widths, same line breaks, same page count. On a machine that has
% the real font (Word, Overleaf) the first branch wins and nothing moves.
\IfFontExistsTF{Times New Roman}{
  \newcommand{\thesisserif}{Times New Roman}
  \newcommand{\thesissans}{Arial}
  \newcommand{\thesismono}{Courier New}
}{
  \newcommand{\thesisserif}{Liberation Serif}
  \newcommand{\thesissans}{Liberation Sans}
  \newcommand{\thesismono}{Liberation Mono}
}
\setmainfont[Ligatures=TeX]{\thesisserif}
\setsansfont[Ligatures=TeX]{\thesissans}
\setmonofont[Scale=0.90,Ligatures=TeX]{\thesismono}

% ---- Page ------------------------------------------------------------
% 2.54 cm on every side — the margin the department template uses.
\usepackage[a4paper,margin=2.54cm]{geometry}
\usepackage{microtype}
\emergencystretch=1.5em          % prefer loose spacing over an overfull line
\usepackage{setspace}
\onehalfspacing                  % Guidelines A.2: line spacing 1.5
\setlength{\parskip}{0pt}
\setlength{\parindent}{1.2em}

\usepackage{csquotes}
\usepackage{enumitem}
\setlist{itemsep=2pt,parsep=0pt,topsep=4pt}

% ---- Tables and figures ---------------------------------------------
\usepackage{booktabs}
\usepackage{tabularx}
\usepackage{array}
\usepackage{multirow}
\usepackage{graphicx}
\usepackage{xcolor}
\usepackage[font=small,labelfont=bf,textfont=small]{caption}
\usepackage{subcaption}
\usepackage{float}

% ---- Diagrams --------------------------------------------------------
\usepackage{tikz}
\usetikzlibrary{positioning,arrows.meta,calc,fit,backgrounds,shapes.geometric,
                decorations.pathreplacing}
\usepackage{pgfplots}
\pgfplotsset{compat=1.18}
\usepgfplotslibrary{groupplots}
% =====================================================================
%  diagrams.tex — THE diagram standard. One file, every figure.
%
%  The rule this file exists to enforce: a reader should be able to look
%  at any two figures in the thesis and see the same visual language.
%  Same palette, same line weights, same type size, same width.
%
%  There are exactly TWO kinds of figure, and every figure is one of them:
%
%    dgdata  — numbers on axes (pgfplots). Anything measured or computed.
%    dgflow  — structure and sequence (TikZ). Anything explained.
%
%  Nothing in content/ may define a colour, a width, or a font size.
%  tools/diagram-check.py fails the build if it does.
% =====================================================================

% ---------------------------------------------------------------------
%  Palette — Okabe & Ito, chosen because it survives the three common
%  colour-vision deficiencies AND greyscale printing. A thesis gets
%  printed in black and white more often than anyone admits.
%
%  Each colour has ONE meaning across the whole thesis:
% ---------------------------------------------------------------------
\definecolor{dgCFT}{HTML}{0072B2}   % crash-fault protocols            (blue)
\definecolor{dgBFT}{HTML}{D55E00}   % Byzantine-fault protocols     (vermilion)
\definecolor{dgAux}{HTML}{009E73}   % the third series / "correct"     (green)
\definecolor{dgSand}{HTML}{E69F00}  % the fourth series / "attention" (orange)
\definecolor{dgPurple}{HTML}{CC79A7}% the fifth series                (purple)
\definecolor{dgGrey}{HTML}{6E6E6E}  % structure only — never a series   (grey)
\definecolor{dgInk}{HTML}{1A1A1A}   % text and axes
\definecolor{dgPaper}{HTML}{FFFFFF} % label backgrounds over lines

% Colour is never the only difference between two series. Every series
% also gets its own mark and its own dash pattern, so the figure still
% reads on a photocopy. These four cycles stay in lockstep.
\pgfplotscreateplotcyclelist{dgcycle}{%
  {dgCFT,   mark=*,        mark size=1.9pt, solid},
  {dgBFT,   mark=square*,  mark size=1.8pt, dashed},
  {dgAux,   mark=triangle*,mark size=2.1pt, dotted},
  {dgSand,  mark=diamond*, mark size=2.1pt, dashdotted},
  {dgPurple,mark=pentagon*,mark size=1.9pt, densely dashed},
}

% ---------------------------------------------------------------------
%  Fixed geometry. Every figure is the same width, so the eye does not
%  have to re-scale between them.
% ---------------------------------------------------------------------
% Computed at \begin{document}, NOT here: in the preamble \linewidth is
% still the class default (360pt), not the real text block set by
% geometry. Setting it here silently made every figure 10.9cm wide on a
% 15.9cm page.
\newlength{\dgwidth}
\newlength{\dgheight}
\AtBeginDocument{%
  \setlength{\dgwidth}{0.86\linewidth}%
  \setlength{\dgheight}{0.50\dgwidth}%
}

% ---------------------------------------------------------------------
%  KIND 1 — dgdata: numbers on axes
%
%    \begin{dgdata}{xlabel={...}, ylabel={...}}
%      \addplot ... ;
%    \end{dgdata}
%
%  Everything that could drift between figures is set here once.
% ---------------------------------------------------------------------
\pgfplotsset{
  dgbase/.style={
    width=\dgwidth, height=\dgheight,
    cycle list name=dgcycle,
    every axis plot/.append style={line width=0.9pt},
    tick label style={font=\footnotesize, color=dgInk},
    label style={font=\small, color=dgInk},
    title style={font=\small\bfseries, color=dgInk},
    legend style={font=\footnotesize, draw=dgGrey!45, fill=dgPaper,
                  fill opacity=0.92, text opacity=1,
                  cells={anchor=west}, row sep=1pt},
    legend cell align=left,
    axis line style={color=dgGrey!70, line width=0.6pt},
    grid=major,
    major grid style={color=dgGrey!22, line width=0.4pt},
    tick align=outside, tick style={color=dgGrey!70},
    enlarge x limits=0.06, enlarge y limits=0.10,
    axis on top,
  },
  % Bar charts share the same skin but need their own spacing.
  dgbars/.style={
    dgbase,
    ybar, bar width=8pt, enlarge x limits=0.14,
    ymin=0, grid=major, ymajorgrids=true, xmajorgrids=false,
    % "fill=." means "fill with this plot's own cycle colour". Without it
    % the cycle list sets only the line colour and every bar comes out as
    % an empty outline.
    every axis plot/.append style={fill=., fill opacity=0.85,
                                   draw=dgInk!55, line width=0.4pt},
    xtick style={draw=none},
  },
}

\newenvironment{dgdata}[1][]
  {\begin{tikzpicture}\begin{axis}[dgbase,#1]}
  {\end{axis}\end{tikzpicture}}

\newenvironment{dgbarchart}[1][]
  {\begin{tikzpicture}\begin{axis}[dgbars,#1]}
  {\end{axis}\end{tikzpicture}}

% ---------------------------------------------------------------------
%  KIND 2 — dgflow: structure and sequence
%
%    \begin{dgflow}
%      \node[dgbox] (a) {...};
%      \node[dgbox, right=of a] (b) {...};
%      \draw[dgarrow] (a) -- (b) node[dgtag,midway]{...};
%    \end{dgflow}
%
%  Node styles are deliberately few. Five box styles is already too many
%  for a reader to hold; these are four, and each means one thing.
% ---------------------------------------------------------------------
\tikzset{
  dgflowbase/.style={
    node distance=9mm and 13mm,
    every node/.append style={font=\small, text=dgInk, align=center},
    >={Stealth[length=2.4mm,width=1.8mm]},
  },
  % a thing that exists: a node, a component, a layer
  dgbox/.style={
    draw=dgGrey!75, line width=0.6pt, rounded corners=1.5pt,
    fill=dgGrey!7, inner sep=4pt, minimum height=7mm, minimum width=20mm,
  },
  % the same, but this one is the subject of the figure
  dgboxkey/.style={dgbox, draw=dgCFT, line width=0.9pt, fill=dgCFT!8},
  % the same, but this one is where it goes wrong
  dgboxbad/.style={dgbox, draw=dgBFT, line width=0.9pt, fill=dgBFT!8},
  % a moment in time, or a state
  dgstate/.style={
    draw=dgGrey!75, line width=0.6pt, circle, fill=dgPaper,
    inner sep=2.5pt, minimum size=9mm,
  },
  % movement: a message, a step, a dependency
  dgarrow/.style={->, draw=dgGrey!85, line width=0.7pt},
  dgarrowkey/.style={dgarrow, draw=dgCFT, line width=0.9pt},
  dgarrowbad/.style={dgarrow, draw=dgBFT, line width=0.9pt, dashed},
  % a label sitting on a line — the white background is what stops it
  % colliding with the line underneath
  % NOTE: text=, not color=. TikZ's color= sets the fill too, so
  % "color=dgInk" here would override fill=dgPaper and render every tag
  % as dark text on a dark box.
  dgtag/.style={
    font=\footnotesize, text=dgInk, fill=dgPaper, inner sep=1.5pt,
  },
  % an aside: the sentence that explains what the reader is looking at
  dgnote/.style={
    font=\footnotesize\itshape, color=dgGrey, align=left, inner sep=2pt,
  },
  % an axis inside a flow figure (a timeline, a spectrum)
  dgaxis/.style={draw=dgGrey!70, line width=0.6pt, ->},
}

\newenvironment{dgflow}[1][]
  {\begin{tikzpicture}[dgflowbase,#1]}
  {\end{tikzpicture}}

% ---------------------------------------------------------------------
%  Figure wrapper. Used by every figure so the spacing above and below
%  a figure is identical everywhere.
%
%    \dgfigure{label}{short caption}{long caption}{ ...body... }
% ---------------------------------------------------------------------
\newcommand{\dgfigure}[4]{%
  \begin{figure}[htbp]
    \centering
    #4
    \caption[#2]{#3}
    \label{#1}
  \end{figure}}


% ---- Maths -----------------------------------------------------------
% amssymb is NOT loaded: unicode-math (below) provides every
% symbol it defines, and loading both makes them fight over
% \eth, \digamma, \smallsetminus and \backepsilon.
\usepackage{amsmath}

% The maths FONT.
% Without this the text is Times (or its metric twin) and the
% mathematics is Computer Modern, because that is what XeLaTeX uses when
% nothing says otherwise. The two do not match, and the mismatch shows
% wherever a symbol sits inside a sentence: "FaB needs $5f+1$ nodes"
% printed the words in one face and the 5f+1 in another.
%
% TeX Gyre Termes Math is the companion of Times: same proportions, same
% weight, designed for exactly this pairing. If it is missing the
% document still builds, in Computer Modern maths, and says so in the
% log rather than failing.
%
% NOT the text font: TeX Gyre Termes has no accented Greek — no ά, έ, ή,
% ί, ό, ύ, ώ — so the Greek edition would be full of missing glyphs.
% Liberation Serif has the full range, which is why it is the fallback.
\usepackage{unicode-math}
\IfFontExistsTF{texgyretermes-math.otf}{
  \setmathfont{texgyretermes-math.otf}
  % unicode-math sets upright Greek capitals by ISO style; LaTeX's own
  % convention, and every source this thesis quotes, prints them upright.
  \setmathfont[range=\mathup/{greek,Greek}]{texgyretermes-math.otf}
}{
  \PackageWarningNoLine{thesis}{TeX Gyre Termes Math not found —
    mathematics will be Computer Modern and will not match the text.
    Install the tex-gyre-math package}
}


% ---- Citations and links ---------------------------------------------
% cite sorts and compresses numeric citations: \cite{a,b,c} -> [2]-[4].
% It must load BEFORE hyperref.
\usepackage{cite}
\usepackage[colorlinks=true,linkcolor=black,citecolor=black,
            urlcolor=dgCFT,breaklinks=true]{hyperref}
\usepackage{url}
\urlstyle{same}

% ---- Source-quality markers ------------------------------------------
% Project rule R9: a source that has not been through peer review is
% marked WHERE IT IS USED, not only in the bibliography. A reader
% checking a claim should not have to turn to the back to find out that
% the evidence is a preprint.
% Capitals, not small capitals: Times New Roman has no true small-caps
% face, and neither does the metric-compatible fallback, so \textsc{}
% only produced a font warning and a synthesised shape.
\newcommand{\grey}{\textsuperscript{G}}   % grey literature
\newcommand{\prep}{\textsuperscript{P}}   % preprint, not refereed
% U marks a DIFFERENT thing from G and P. Those two say what kind of
% publication a source is. U says that the thesis could not verify the
% source's own bibliographic data against the article or an index, and
% took it from another work's reference list. A reader who wants to
% follow the citation needs to know that before they try.
\newcommand{\unver}{\textsuperscript{U}}   % bibliographic data unverified

% ---- Headings --------------------------------------------------------
% Guidelines A.2: chapter headings 14 pt bold with 12 pt before and
% after; sub-headings 12 pt bold; every chapter on a new page. The
% template sets a chapter as two lines: "ΚΕΦΑΛΑΙΟ 1" flush left, the
% title centred under it. docx/make-reference.py gives the .docx the
% same (Chapter Label: left; Heading 1: centred, 14 pt).
\usepackage{titlesec}
\titleformat{\chapter}[display]
  {\bfseries\fontsize{14}{17}\selectfont}
  {\filright\chaptertitlename\ \thechapter}{6pt}{\filcenter}
\titlespacing*{\chapter}{0pt}{12pt}{12pt}
\titleformat{\section}
  {\bfseries\fontsize{12}{15}\selectfont}{\thesection}{0.8em}{}
\titleformat{\subsection}
  {\bfseries\fontsize{12}{15}\selectfont}{\thesubsection}{0.8em}{}

% ---- Front-matter heading -------------------------------------------
% 14 pt bold, centred — the department template's look for the pages
% before Chapter 1. It is an unnumbered \chapter* (set by titlesec
% above), so these headings sit at the same height as TABLE OF CONTENTS,
% the lists and REFERENCES.
%   #1 = as printed   #2 = as it appears in the table of contents
\newcommand{\frontheading}[2]{%
  \chapter*{#1}%
  \phantomsection
  \addcontentsline{toc}{chapter}{#2}%
  \markboth{#2}{#2}}

% Blank leaf, as the template has between front-matter sections. It is
% counted and numbered, so both editions and the .docx agree page for page.
\newcommand{\blankleaf}{%
  \clearpage
  \thispagestyle{plain}%
  \null\vfill\begin{center}\small\dgblankleaftext\end{center}\vfill\null
  \clearpage}

% ---- Measurements that do not exist yet ------------------------------
% Project rule R1: never print a number that was not measured. Every
% figure from the author's own benchmark campaign is declared once in
% data/measurements.tex. Until a value is filled in, \meas prints a
% visible marker instead of a plausible-looking number, and
% tools/data-check.py lists what is still outstanding.
\newcommand{\PENDING}{\textbf{\color{dgBFT}[PENDING]}}
\newcommand{\meas}[1]{\csname measvalue@#1\endcsname}
\newcommand{\declaremeasurement}[2]{%
  \expandafter\newcommand\csname measvalue@#1\endcsname{#2}}
\newcommand{\pendingmeasurement}[1]{%
  \expandafter\newcommand\csname measvalue@#1\endcsname{\PENDING}}

% =====================================================================
%  lang-el.tex — η ελληνική έκδοση.
%
%  Η ελληνική ΑΚΟΛΟΥΘΕΙ την αγγλική: το content/en/ είναι το πρωτότυπο,
%  το content/el/ η μετάφρασή του. Το tools/sync-check.py δείχνει ποια
%  ελληνικά κομμάτια έχουν μείνει πίσω. Δες docs/SYNC.md.
% =====================================================================
\setmainlanguage[numerals=arabic]{greek}
\setotherlanguage{english}

% Ελληνικές γραμματοσειρές: ίδιες με τις λατινικές, ώστε οι δύο εκδόσεις
% να μοιάζουν με το ίδιο έγγραφο.
\newfontfamily\greekfont[Ligatures=TeX]{\thesisserif}
\newfontfamily\greekfontsf[Ligatures=TeX]{\thesissans}
\newfontfamily\greekfonttt[Scale=0.90,Ligatures=TeX]{\thesismono}

% \en{...} : αγγλικός όρος μέσα σε ελληνικό κείμενο, με σωστό συλλαβισμό
\newcommand{\en}[1]{\textenglish{\emph{#1}}}
% \enup{...}: αγγλικό όνομα χωρίς πλάγια (Paxos, HotStuff, …)
\newcommand{\enup}[1]{\textenglish{#1}}
\newcommand{\gr}[1]{#1}

% Ελληνική σύμβαση αριθμών στα σχήματα: υποδιαστολή κόμμα.
\pgfkeys{/pgf/number format/.cd, use comma, 1000 sep={.}}

\newcommand{\dgblankleaftext}{(Κενό φύλλο)}

\newcommand{\thesisnames}{%
  \renewcommand{\chaptername}{ΚΕΦΑΛΑΙΟ}%
  \renewcommand{\appendixname}{ΠΑΡΑΡΤΗΜΑ}%
  \renewcommand{\figurename}{Σχήμα}%
  \renewcommand{\tablename}{Πίνακας}%
  \renewcommand{\bibname}{ΒΙΒΛΙΟΓΡΑΦΙΑ}%
  \renewcommand{\contentsname}{ΠΕΡΙΕΧΟΜΕΝΑ}%
  \renewcommand{\listfigurename}{ΚΑΤΑΛΟΓΟΣ ΣΧΗΜΑΤΩΝ}%
  \renewcommand{\listtablename}{ΚΑΤΑΛΟΓΟΣ ΠΙΝΑΚΩΝ}%
}

% ΕΛΛΗΝΙΚΕΣ ετικέτες σχημάτων. Το labels-en.tex είναι το αγγλικό δίδυμο: τα δύο
% αρχεία πρέπει να ορίζουν ΑΚΡΙΒΩΣ το ίδιο σύνολο κλειδιών, με την ίδια σειρά —
% ένα κλειδί που υπάρχει στο ένα και όχι στο άλλο τυπώνει σιωπηλό κενό.
%
% Οι πηγές των σχημάτων δεν περιέχουν δικό τους αναγνώσιμο κείμενο: κάθε
% συμβολοσειρά είναι \fl{key} και κάθε κλειδί ορίζεται εδώ. Ονόματα
% πρωτοκόλλων, συνεδρίων, χρονολογίες, αριθμοί, μαθηματικά και μονάδες είναι
% ίδια στις δύο εκδόσεις, μένουν αυτούσια στον κώδικα και ΔΕΝ είναι κλειδιά.

\providecommand{\deflabel}[2]{\expandafter\newcommand\csname figlabel@#1\endcsname{#2}}
\providecommand{\fl}[1]{\csname figlabel@#1\endcsname}

\deflabel{equivocation.byzantine.note}{ο $K_2$ ζει και συνεχίζει\\\textbf{αντιφατική δήλωση}}
\deflabel{equivocation.byzantine.title}{ΒΥΖΑΝΤΙΝΟ}
\deflabel{equivocation.crash.note}{ο $K_2$ σταματά\\ίδιο μήνυμα σε όλους}
\deflabel{equivocation.crash.title}{ΚΑΤΑΡΡΕΥΣΗ}
\deflabel{quorum.byzantine}{βυζαντινό, κανονική περίπτωση}
\deflabel{quorum.crash}{κατάρρευση (και βυζαντινό με έμπιστο υλικό)}
\deflabel{quorum.fairorder}{δίκαιη διάταξη}
\deflabel{quorum.onephaseless}{μία φάση λιγότερη}
\deflabel{quorum.twostep}{κάτω φράγμα δύο βημάτων}
\deflabel{quorum.xlabel}{πλήθος ανεκτών σφαλμάτων}
\deflabel{quorum.ylabel}{απαιτούμενοι κόμβοι}
\deflabel{quote.close}{\guillemotright}
\deflabel{quote.open}{\guillemotleft}
\deflabel{scaling.legend.geo}{γεωκατανεμημένο}
\deflabel{scaling.panel.geo}{(β) LAN έναντι γεωκατανεμημένου}
\deflabel{scaling.panel.lan}{(α) ένα \en{datacenter} (LAN)}
\deflabel{scaling.ylabel}{μείωση \en{throughput}}
\deflabel{smr.agreement.order}{συμφωνία $+$ διάταξη}
\deflabel{smr.clients}{πελάτες}
\deflabel{smr.consensus}{ΣΥΝΑΙΝΕΣΗ /\\ΑΤΟΜΙΚΗ ΕΚΠΟΜΠΗ}
\deflabel{smr.deterministic}{ΝΤΕΤΕΡΜΙΝΙΣΤΙΚΕΣ}
\deflabel{smr.here.live}{εδώ ζουν}
\deflabel{smr.machine}{μηχανή}
\deflabel{smr.same.io}{ίδια είσοδος $\Rightarrow$ ίδια έξοδος}
\deflabel{sync.liveness.guaranteed}{ζωτικότητα εγγυημένη}
\deflabel{sync.liveness.none}{ζωτικότητα: καμία εγγύηση}
\deflabel{sync.model.asynchronous}{Ασύγχρονο}
\deflabel{sync.model.partial}{Μερικώς σύγχρονο}
\deflabel{sync.model.synchronous}{Σύγχρονο}
\deflabel{sync.note.asynchronous}{κανένα φράγμα}
\deflabel{sync.note.safety}{Η ασφάλεια ισχύει σε ΟΛΟ τον άξονα, και πριν και μετά την GST.}
\deflabel{sync.note.synchronous}{κάθε μήνυμα $\le\Delta$, και το $\Delta$ είναι γνωστό}
\deflabel{timeline.milestone.byzmodel}{βυζαντινό μοντέλο σφαλμάτων}
\deflabel{timeline.milestone.clocks}{λογικά ρολόγια}
\deflabel{timeline.milestone.linearizability}{γραμμικοποιησιμότητα}
\deflabel{timeline.milestone.partialsync}{μερική συγχρονία}
\deflabel{timeline.milestones.title}{Θεωρητικά ορόσημα (δεν είναι πρωτόκολλα):}

% =====================================================================
%  meta.tex — the facts about this thesis, written once.
%
%  Both editions and both .docx files read this file. A title or a name
%  changed here changes everywhere; a title changed anywhere else changes
%  in one place and silently disagrees with the other three.
%
%  docx/frontmatter.py parses this file directly. Keep the
%  \newcommand{\th...}{...} shape.
% =====================================================================
\newcommand{\thTitleEN}{Evolution and Comparison of Consensus Protocols
                        in Distributed Systems}
\newcommand{\thTitleEL}{Εξέλιξη και Σύγκριση Πρωτοκόλλων Συναίνεσης
                        σε Κατανεμημένα Συστήματα}

\newcommand{\thAuthorEN}{Panagiotis Dampanis}
\newcommand{\thAuthorEL}{Παναγιώτης Δαμπάνης}
\newcommand{\thAM}{711070095}

\newcommand{\thSupervisorEN}{Grammati Pantziou, Professor}
\newcommand{\thSupervisorEL}{Γραμματή Παντζιού, Καθηγήτρια}

\newcommand{\thUniversityEN}{UNIVERSITY OF WEST ATTICA}
\newcommand{\thUniversityEL}{ΠΑΝΕΠΙΣΤΗΜΙΟ ΔΥΤΙΚΗΣ ΑΤΤΙΚΗΣ}

% English names as the institution prints them (checked 2026-11-11):
% "Department of Informatics and Computer Engineering" on ice.uniwa.gr/en/;
% "School of Engineering" above that department on
% uniwa.gr/en/studies/schools-and-departments/. Capitals, like the
% Greek lines and \thUniversityEN.
\newcommand{\thSchoolEN}{SCHOOL OF ENGINEERING}
\newcommand{\thSchoolEL}{ΣΧΟΛΗ ΜΗΧΑΝΙΚΩΝ}
\newcommand{\thDepartmentEN}{DEPARTMENT OF INFORMATICS AND COMPUTER ENGINEERING}
\newcommand{\thDepartmentEL}{ΤΜΗΜΑ ΜΗΧΑΝΙΚΩΝ ΠΛΗΡΟΦΟΡΙΚΗΣ ΚΑΙ ΥΠΟΛΟΓΙΣΤΩΝ}

% Examination committee — filled in when the committee is appointed.
% Two members besides the supervisor, who is listed on her own line:
% that is the template's title page and both of the advisor's recent
% theses (2025, 2026).
\newcommand{\thExamA}{\thTODO{NAME, SURNAME, RANK}}
\newcommand{\thExamB}{\thTODO{NAME, SURNAME, RANK}}

\newcommand{\thBlank}[1]{\rule{#1}{0.4pt}}
\newcommand{\thExamDateEN}{Date of examination
  \thBlank{1.2cm}/\thBlank{1.2cm}/\thBlank{2cm}}
\newcommand{\thExamDateEL}{Ημερομηνία εξέτασης
  \thBlank{1.2cm}/\thBlank{1.2cm}/\thBlank{2cm}}

% Anything not yet known is marked, not guessed. tools/todo-check.py
% counts these and the build prints the count, so a placeholder cannot
% reach the examiners unnoticed.
\newcommand{\thTODO}[1]{\textbf{\color{dgBFT}$\langle$#1$\rangle$}}

\newcommand{\thScientificAreaEN}{Distributed systems --- consensus
  protocols and performance evaluation}
\newcommand{\thScientificAreaEL}{Κατανεμημένα συστήματα ---
  πρωτόκολλα συναίνεσης και αποτίμηση επιδόσεων}

% =====================================================================
%  measurements.tex — every number that must come from YOUR OWN
%  benchmark campaign, declared in one place.
%
%  WHY THIS FILE EXISTS
%
%  Project rule R1: never print a number that was not measured. The way
%  that rule gets broken is never on purpose. It gets broken because a
%  plausible number goes into a draft as a placeholder, the draft gets
%  edited forty times, and by the submission the number looks like it has
%  always been there.
%
%  So a number that does not exist yet cannot be typed into the text at
%  all. It is declared here as pending, the text asks for it by name with
%  \meas{key}, and until you fill it in the page prints a red [PENDING]
%  that no examiner could miss and no author could forget.
%
%  HOW TO FILL ONE IN
%
%    \pendingmeasurement{raft.throughput.n3}          <- before
%    \declaremeasurement{raft.throughput.n3}{12\,400} <- after
%
%  Then run `make data` — it lists what is still outstanding.
%
%  NAMING
%
%    <system>.<quantity>.<condition>
%    etcd.p99.n5.open      etcd, 99th percentile latency, 5 nodes, open loop
%
%  Everything in this file traces back to a research question in §1.4 and
%  to the four comparison criteria named there: consensus latency,
%  throughput, fault tolerance, scaling.
% =====================================================================

% ---------------------------------------------------------------------
%  RQ4 — empirical behaviour of selected production implementations
%  (§1.4, "ΕΕ4 / RQ4"). Small controlled cluster.
%
%  Instrument: one load generator for every system, open loop for
%  throughput and latency tails, closed loop for the saturation search
%  (stated in the abstract). Values below are per implementation, at the
%  cluster sizes the budget allows.
% ---------------------------------------------------------------------

% -- Throughput at saturation, requests/s ------------------------------
\pendingmeasurement{etcd.tput.n3}         % etcd (Raft), 3 nodes
\pendingmeasurement{etcd.tput.n5}         % etcd (Raft), 5 nodes
\pendingmeasurement{etcd.tput.n7}         % etcd (Raft), 7 nodes
\pendingmeasurement{tendermint.tput.n4}   % Tendermint, 4 nodes  (3f+1, f=1)
\pendingmeasurement{tendermint.tput.n7}   % Tendermint, 7 nodes  (3f+1, f=2)
\pendingmeasurement{hotstuff.tput.n4}     % HotStuff, 4 nodes
\pendingmeasurement{hotstuff.tput.n7}     % HotStuff, 7 nodes

% -- Latency, milliseconds. Median AND tail: the tail is the point ------
\pendingmeasurement{etcd.p50.n3}
\pendingmeasurement{etcd.p99.n3}
\pendingmeasurement{etcd.p50.n5}
\pendingmeasurement{etcd.p99.n5}
\pendingmeasurement{tendermint.p50.n4}
\pendingmeasurement{tendermint.p99.n4}
\pendingmeasurement{hotstuff.p50.n4}
\pendingmeasurement{hotstuff.p99.n4}

% -- Fault tolerance: what a single leader failure costs ---------------
%  Time from the failure to the first committed request afterwards, and
%  how much throughput was lost in between. This is the measurement that
%  turns "view change costs O(n)" from a formula into a number.
\pendingmeasurement{etcd.failover.ms}
\pendingmeasurement{tendermint.failover.ms}
\pendingmeasurement{hotstuff.failover.ms}
\pendingmeasurement{etcd.failover.losspct}
\pendingmeasurement{tendermint.failover.losspct}
\pendingmeasurement{hotstuff.failover.losspct}

% ---------------------------------------------------------------------
%  RQ5 — scaling, experimental half only
%  (§1.4 splits RQ5 deliberately: measured for small clusters, from the
%  literature for large ones. Only the measured half belongs here.)
%
%  The claim under test is the arithmetic of Figure "fig:quorum": if the
%  node count rises from 3 to 7, what actually happens to throughput?
% ---------------------------------------------------------------------
\pendingmeasurement{scaling.etcd.dropn3n7pct}       % % throughput lost, 3 -> 7
\pendingmeasurement{scaling.tendermint.dropn4n7pct}
\pendingmeasurement{scaling.hotstuff.dropn4n7pct}

% ---------------------------------------------------------------------
%  RQ3 — comparability
%
%  Not a performance number: the evidence that the comparison is fair.
%  Same hardware, same load model, same instrument, and a run-to-run
%  spread small enough that the differences reported are larger than the
%  noise. Without this, every other number above is decoration.
% ---------------------------------------------------------------------
\pendingmeasurement{env.vm.spec}          % e.g. "4 vCPU, 8 GB RAM, ..."
\pendingmeasurement{env.network.rtt}      % measured RTT between nodes
\pendingmeasurement{env.repeats}          % runs per configuration
\pendingmeasurement{env.cv.max}           % worst coefficient of variation

% ---------------------------------------------------------------------
%  Predictions, registered BEFORE the measurements
%
%  The abstract promises this: "Predictions are registered before the
%  measurements are taken." A prediction written after the fact is not a
%  prediction. Fill these in when the campaign is designed, not when it
%  is finished — git history is what makes the claim checkable, so commit
%  them in their own commit, before the first result lands.
% ---------------------------------------------------------------------
\pendingmeasurement{pred.tput.ranking}    % predicted order, fastest first
\pendingmeasurement{pred.scaling.shape}   % linear? knee? where?
\pendingmeasurement{pred.failover.ranking}

% =====================================================================
%  abbreviations.tex — one list, both languages.
%
%  RULE (project rule R1 applied to definitions): an expansion goes in
%  this list only if it can be traced to the thesis text itself or to
%  the title of a source the thesis already cites. The third argument of
%  each entry records where. An expansion written from memory is exactly
%  the kind of small unsourced claim an examiner enjoys finding.
%
%  Each language file defines \abbrentry to pick the column it wants.
%    #1 short   #2 English   #3 Greek   #4 provenance (never printed)
% =====================================================================
\newcommand{\abbrlist}{%
  \abbrentry{FLP}
    {Fischer--Lynch--Paterson (impossibility of deterministic consensus
     in an asynchronous system)}
    {Fischer--Lynch--Paterson (αδυναμία ντετερμινιστικής συναίνεσης σε
     ασύγχρονο σύστημα)}
    {§1.1 names Fischer, Lynch and Paterson; used as FLP in §2.4}
  \abbrentry{GST}
    {Global Stabilization Time}
    {Global Stabilization Time --- καθολικός χρόνος σταθεροποίησης}
    {§2.1 spells the term out and gives the abbreviation}
  \abbrentry{PBFT}
    {Practical Byzantine Fault Tolerance}
    {Practical Byzantine Fault Tolerance}
    {title of \cite{castro1999pbft}}
  \abbrentry{SMR}
    {State Machine Replication}
    {State Machine Replication --- αναπαραγωγή μηχανής καταστάσεων}
    {§2.3 spells the term out and gives the abbreviation}
  \abbrentry{VR}
    {Viewstamped Replication}
    {Viewstamped Replication}
    {title of \cite{oki1988viewstamped}; §1.2}
}

% =====================================================================
%  NOT IN THE LIST, AND WHY
%
%  These five appear in the thesis but have no expansion that traces to
%  a source we hold. tools/abbr-check.py reports them on every run so
%  they cannot be quietly forgotten.
%
%    BFT, CFT  appear only as labels inside Figure 1.1, which is
%              rendered as an image — a reader cannot find them in the
%              text at all
%    EPaxos    \cite{moraru2013epaxos} is titled "There Is More Consensus
%              in Egalitarian Parliaments"; it does not contain
%              "Egalitarian Paxos"
%    ZAB       \cite{junqueira2011zab} is titled "Zab: High-Performance
%              Broadcast for Primary-Backup Systems"; and that entry's
%              own verification field records that the paper was not read
%    KIP       no cited source expands the acronym
%
%  The fix is in the text, not in this file: the Guidelines require that
%  a term be introduced in both languages at first use. Introduce them
%  and they can move up into \abbrlist.
% =====================================================================
\newcommand{\abbrpending}{BFT, CFT, EPaxos, ZAB, KIP}

\newcommand{\abbrentry}[4]{}

\begin{document}
\thesisnames

% --- Ενότητα 1η: ΕΙΣΑΓΩΓΙΚΑ ΣΤΟΙΧΕΙΑ ---------------------------------
\pagenumbering{roman}
\setcounter{page}{1}

% Εξώφυλλο — η διάταξη του Προτύπου ΤΜΠΥ.
\begin{titlepage}
\centering
\thispagestyle{empty}

\includegraphics[width=2.6cm]{assets/pada-logo.png}

\vspace{0.8em}
{\large \thUniversityEL\par}
{\large \thSchoolEL\par}
{\large \thDepartmentEL\par}

\vfill
{\Large\bfseries ΔΙΠΛΩΜΑΤΙΚΗ ΕΡΓΑΣΙΑ\par}
\vspace{1.5em}
{\LARGE\bfseries \thTitleEL\par}
\vspace{0.8em}
{\large\enup{\itshape \thTitleEN}\par}
\vfill
{\Large\bfseries \thAuthorEL\par}
{\large Α.Μ. \thAM\par}
\vspace{2em}
{\large Εισηγήτρια: \thSupervisorEL\par}
\vfill
\end{titlepage}

% Το περιβάλλον titlepage μηδενίζει τον μετρητή σελίδων στο κλείσιμό του.
% Χωρίς αυτό, εξώφυλλο και επόμενη σελίδα είναι και οι δύο «σελίδα 1» και
% όλη η λατινική αρίθμηση βγαίνει μία πίσω από εκείνη του .docx.
\setcounter{page}{2}

\blankleaf
% Σελίδα τίτλου — επαναλαμβάνει το εξώφυλλο και προσθέτει την επιτροπή.
\thispagestyle{plain}
\begin{center}
\vspace*{1cm}
{\Large\bfseries ΔΙΠΛΩΜΑΤΙΚΗ ΕΡΓΑΣΙΑ\par}
\vspace{1.5em}
{\LARGE\bfseries \thTitleEL\par}
\vspace{0.8em}
{\large\enup{\itshape \thTitleEN}\par}
\vspace{2.5em}
{\Large\bfseries \thAuthorEL\par}
{\large Α.Μ. \thAM\par}
\vspace{2.5em}
{\large Εισηγήτρια:\par}
{\large \thSupervisorEL\par}
\vspace{1.4cm}% room for the digital signature, as in the advisor's theses
{\large Εξεταστική Επιτροπή:\par}
\vspace{0.4em}
{\large \thExamA\par}
\vspace{1.4cm}
{\large \thExamB\par}
\vspace{1.4cm}
{\large \thExamDateEL\par}
\end{center}
\clearpage

\blankleaf
% =====================================================================
%  ΔΗΛΩΣΗ ΣΥΓΓΡΑΦΕΑ
%
%  ΤΟ ΚΕΙΜΕΝΟ ΕΙΝΑΙ ΑΥΤΟΥΣΙΟ από το «Πρότυπο Συγγραφής Διπλωματικής
%  Εργασίας» του ΤΜΠΥ (PROTYPO-SYGGRAFIS-TMPY.docx, παρ. 188-200).
%  ΜΗΝ ΤΟ ΠΕΙΡΑΞΕΙΣ. Δεν το συντάσσει ο φοιτητής.
%
%  ΤΙ ΔΕΝ ΔΙΟΡΘΩΘΗΚΕ ΣΚΟΠΙΜΑ, γιατί έτσι είναι στο Πρότυπο:
%   * ο τίτλος λέει «ΜΕΤΑΠΤΥΧΙΑΚΗΣ» ενώ η εργασία είναι διπλωματική.
%  Αν το Τμήμα δώσει διορθωμένη έκδοση, αντικατέστησε ΟΛΟ το αρχείο.
%
%  ΤΙ ΔΙΑΦΕΡΕΙ ΑΝΑΓΚΑΣΤΙΚΑ από το Πρότυπο (τυπογραφία, όχι κείμενο):
%   * τα διπλά κενά του Προτύπου γίνονται μονά — το TeX ενώνει τα κενά
%     και το ίδιο κάνει και ο αναλυτής του DOCX·
%   * οι τελείες στη γραμμή «του …………» τις βάζει το \dotfill·
%   * κάτω από το «Ο/Η Δηλών/ούσα» μπαίνει το όνομα, ενώ το Πρότυπο
%     αφήνει κενό χώρο για υπογραφή. Σβήσ' το αν υπογράφεις με το χέρι.
%
%  ΤΟ ΙΔΙΟ ΑΡΧΕΙΟ διαβάζει και το docx/frontmatter.py. Η αγγλική
%  έκδοση έχει ΜΕΤΑΦΡΑΣΗ ΕΥΚΟΛΙΑΣ (frontmatter/en/declaration.tex),
%  που το δηλώνει πάνω στη σελίδα: δεσμευτικό είναι ΑΥΤΟ το κείμενο.
% =====================================================================
\frontheading{ΔΗΛΩΣΗ ΣΥΓΓΡΑΦΕΑ ΜΕΤΑΠΤΥΧΙΑΚΗΣ ΕΡΓΑΣΙΑΣ}{Δήλωση Συγγραφέα}

Ο/η κάτωθι υπογεγραμμένος/η \thAuthorEL{} του \dotfill{}, με αριθμό
μητρώου \thAM{} φοιτητής/τρια του Τμήματος Μηχανικών Πληροφορικής και
Υπολογιστών της Σχολής Μηχανικών του Πανεπιστημίου Δυτικής Αττικής,
δηλώνω ότι:

«Βεβαιώνω ότι είμαι συγγραφέας αυτής της Διπλωματικής εργασίας και κάθε
βοήθεια την οποία είχα για την προετοιμασία της, είναι πλήρως
αναγνωρισμένη και αναφέρεται στην εργασία. Επίσης, οι όποιες πηγές από
τις οποίες έκανα χρήση δεδομένων, ιδεών ή λέξεων, είτε ακριβώς είτε
παραφρασμένες, αναφέρονται στο σύνολό τους, με πλήρη αναφορά στους
συγγραφείς, τον εκδοτικό οίκο ή το περιοδικό, συμπεριλαμβανομένων και
των πηγών που ενδεχομένως χρησιμοποιήθηκαν από το διαδίκτυο. Επίσης,
βεβαιώνω ότι αυτή η εργασία έχει συγγραφεί από μένα αποκλειστικά και
αποτελεί προϊόν πνευματικής ιδιοκτησίας τόσο δικής μου, όσο και του
Ιδρύματος.

Παράβαση της ανωτέρω ακαδημαϊκής μου ευθύνης αποτελεί ουσιώδη λόγο για
την ανάκληση του πτυχίου μου».

«Βεβαιώνω ότι είμαι συγγραφέας της παρούσας διπλωματικής εργασίας και
ότι έχω αναφέρει ή παραπέμψει σε αυτή, ρητά και συγκεκριμένα, όλες τις
πηγές από τις οποίες έκανα χρήση δεδομένων, ιδεών, προτάσεων ή λέξεων,
είτε αυτές μεταφέρονται επακριβώς (στο πρωτότυπο ή μεταφρασμένες) είτε
παραφρασμένες. Επίσης βεβαιώνω ότι αυτή η εργασία προετοιμάστηκε από
εμένα προσωπικά ειδικά για την συγκεκριμένη διπλωματική εργασία»

\begin{center}
Ο/Η Δηλών/ούσα
\end{center}

\begin{center}
\thAuthorEL
\end{center}

\blankleaf
% =====================================================================
%  ΕΥΧΑΡΙΣΤΙΕΣ — προαιρετική σελίδα, και η μόνη της εργασίας που είναι
%  αποκλειστικά δική σου. Το παρακάτω είναι προσχέδιο· ξαναγράψ' το με
%  δικά σου λόγια.
%
%  Τρεις συμβάσεις που τηρεί:
%    1. πρώτα το ακαδημαϊκό, μετά το προσωπικό
%    2. λέει ΤΙ έκανε ο καθένας, όχι πόσο σπουδαίος είναι
%    3. καμία ιδιωτική πληροφορία — η εργασία δημοσιεύεται και ευρετηριάζεται
% =====================================================================
\frontheading{ΕΥΧΑΡΙΣΤΙΕΣ}{Ευχαριστίες}

Ευχαριστώ την επιβλέπουσα καθηγήτρια Γραμματή Πάντζιου, που δέχτηκε ένα
θέμα με μεγαλύτερο πειραματικό σκέλος από το συνηθισμένο. Μου άφησε τον
χώρο να το δουλέψω με τη δική μου μεθοδολογία, και δεν έπαψε να θέτει τα
σωστά ερωτήματα στα σημεία που είχε σημασία.

Ευχαριστώ τους γονείς μου, που στήριξαν μια φοιτητική διαδρομή αρκετά πιο
μακριά από την προγραμματισμένη και δεν με ρώτησαν ούτε μία φορά πότε θα
τελειώσει. Και την αδελφή μου, που ήξερε πότε να ρωτήσει και πότε όχι.

Η εργασία γράφτηκε παράλληλα με πλήρη απασχόληση. Ό,τι λάθος έμεινε μέσα
είναι δικό μου.

\vspace{2em}
\begin{flushright}
\thAuthorEL\\
Αθήνα, \the\year
\end{flushright}

\blankleaf
% =====================================================================
%  ΠΕΡΙΛΗΨΗ — μετάφραση του frontmatter/en/abstract.tex.
%
%  ΚΑΝΟΝΑΣ: η περίληψη γράφεται ΤΕΛΕΥΤΑΙΑ. Ό,τι υπόσχεται πρέπει να
%  υπάρχει στο κυρίως κείμενο. Το σημείο [ΣΥΜΠΛΗΡΩΣΕ] είναι εκεί όπου
%  μπαίνουν τα ευρήματα, και μένει ώσπου να τρέξει η καμπάνια μετρήσεων.
%  Κανόνας R1: κανένας αριθμός που δεν μετρήθηκε.
%
%  Οδηγίες ΤΜΠΥ: 250-300 λέξεις, πέντε λέξεις-κλειδιά. Το «make check»
%  τα μετράει και τα δύο.
% =====================================================================
\frontheading{ΠΕΡΙΛΗΨΗ}{Περίληψη}

Η συναίνεση είναι ο μηχανισμός με τον οποίο ένα κατανεμημένο σύστημα
συμφωνεί σε μία κοινή σειρά εντολών παρά τα σφάλματα και τις
καθυστερήσεις. Από το 1988 ως σήμερα έχουν προταθεί δεκάδες πρωτόκολλα,
και η βιβλιογραφία δεν έχει καταλήξει σε ενιαία βάση σύγκρισής τους. Οι
δημοσιευμένες μετρήσεις δεν μεταφέρονται μεταξύ εργασιών, επειδή
εξαρτώνται από το υλικό, το μοντέλο φόρτου και την τοπολογία ανάπτυξης. Η
ίδια σύγκριση αντιστρέφεται επίσης ανάλογα με το ποια μετρική επιλέγεται.

Η εργασία μελετά την εξέλιξη των πρωτοκόλλων συναίνεσης και συγκρίνει
επιλεγμένες \emph{υλοποιήσεις} τους σε ελεγχόμενο περιβάλλον. Καλύπτονται
τα πρωτόκολλα κατάρρευσης Paxos, Raft και EPaxos, και τα βυζαντινά
Tendermint και HotStuff, μαζί με τους προδρόμους και τους απογόνους τους.
Κάθε πρωτόκολλο περιγράφεται από την πρωτογενή του πηγή. Κάθε μετακίνηση
στον χώρο σχεδίασης αποδίδεται στη συγκεκριμένη ανταλλαγή που την εξηγεί.

Η αποτίμηση γίνεται με ενιαίο όργανο μέτρησης πάνω σε συστάδα εικονικών
μηχανών, με υβριδικό μοντέλο φόρτου. Τρέχει ανοιχτός βρόχος για τη
ρυθμαπόδοση και τις ουρές καθυστέρησης, και κλειστός για την αναζήτηση
κορεσμού. Οι προβλέψεις δηλώνονται πριν από τις μετρήσεις. Κάθε
συγκριτικός ισχυρισμός φέρει το επίπεδο συγκρισιμότητάς του, ώστε να
διακρίνεται τι αποδίδεται στο πρωτόκολλο και τι στη μηχανική της
υλοποίησης.

\textbf{[ΣΥΜΠΛΗΡΩΣΕ ΜΕΤΑ ΤΑ ΔΕΔΟΜΕΝΑ: τα ευρήματα, σε 2--3 προτάσεις.]}

Η εργασία δηλώνει ρητά και τα όριά της. Το αντικείμενο της μέτρησης είναι
πάντοτε μια υλοποίηση με συγκεκριμένη έκδοση και ρύθμιση, όχι το
πρωτόκολλο καθαυτό. Το εύρος των μεγεθών συστάδας που εξετάζονται
πειραματικά είναι μικρό σε σχέση με τις κλίμακες όπου εμφανίζεται η
ασυμπτωτική συμπεριφορά.

\vspace{1.2em}
\noindent\textbf{ΕΠΙΣΤΗΜΟΝΙΚΗ ΠΕΡΙΟΧΗ:} \thScientificAreaEL.

\vspace{0.6em}
\noindent\textbf{ΛΕΞΕΙΣ ΚΛΕΙΔΙΑ:} πρωτόκολλα συναίνεσης· αναπαραγωγή
μηχανής καταστάσεων· βυζαντινή ανοχή σφαλμάτων· αποτίμηση επιδόσεων·
συγκρισιμότητα.


% Οι Οδηγίες ζητούν την περίληψη ΚΑΙ στις δύο γλώσσες. Το κείμενο είναι
% το ίδιο με της αγγλικής έκδοσης — μία πηγή, που διαβάζεται εδώ αντί να
% αντιγραφεί, ώστε να μην μπορούν να αποκλίνουν.
\cleardoublepage
\begin{english}
% =====================================================================
%  ABSTRACT
%
%  RULE: the abstract is written LAST. Everything it promises must exist
%  in the body. The [PENDING] marker below is where the findings go, and
%  it stays until the benchmark campaign has produced them. Project rule
%  R1: no number that was not measured.
%
%  Guidelines: 250-300 words, five keywords. `make check` measures both
%  and tells you where you stand.
% =====================================================================
\frontheading{ABSTRACT}{Abstract}

Consensus is the mechanism by which a distributed system agrees on a
single order of commands despite faults and delays. Dozens of protocols
have been proposed since 1988, and the literature has not converged on a
common basis for comparing them. Published measurements do not transfer
between papers, because they depend on the hardware, the load model and
the deployment topology. The same comparison also reverses depending on
which metric is chosen.

This thesis studies the evolution of consensus protocols and compares
selected \emph{implementations} of them in a controlled environment. It
covers the crash-fault protocols Paxos, Raft and EPaxos, and the
Byzantine protocols Tendermint and HotStuff, together with their
predecessors and descendants. Each protocol is described from its
primary source. Each move through the design space is attributed to the
specific trade-off that explains it.

The evaluation uses a single measurement instrument on a cluster of
virtual machines, with a hybrid load model. It runs open loop for
throughput and latency tails, and closed loop for the saturation search.
Predictions are registered before the measurements are taken. Every
comparative claim carries its level of comparability, so that what is
attributable to the protocol is separated from what is attributable to
implementation engineering.

\textbf{[FILL IN AFTER THE DATA: the findings, in 2--3 sentences.]}

The limits of the work are stated explicitly. What is measured is always
an implementation at a named version and configuration, not the protocol
itself. The range of cluster sizes examined experimentally is small
relative to the scales at which asymptotic behaviour appears.

\vspace{1.2em}
\noindent\textbf{SCIENTIFIC AREA:} \thScientificAreaEN.

\vspace{0.6em}
\noindent\textbf{KEYWORDS:} consensus protocols; state machine
replication; Byzantine fault tolerance; performance evaluation;
comparability.

\end{english}

\cleardoublepage
\phantomsection
\addcontentsline{toc}{chapter}{Περιεχόμενα}
\tableofcontents

% Σελίδες που απαιτεί το Πρότυπο ΤΜΠΥ ανάμεσα στα περιεχόμενα και το
% Κεφάλαιο 1.

\frontheading{ΚΑΤΑΛΟΓΟΣ ΕΙΚΟΝΩΝ}{Κατάλογος εικόνων}
Η εργασία δεν περιέχει φωτογραφίες ή στιγμιότυπα οθόνης. Όλα τα οπτικά
στοιχεία είναι σχήματα και παρατίθενται στον επόμενο κατάλογο.

\clearpage
\phantomsection
\addcontentsline{toc}{chapter}{Κατάλογος σχημάτων}
\listoffigures

\clearpage
\phantomsection
\addcontentsline{toc}{chapter}{Κατάλογος πινάκων}
\listoftables

\clearpage
\frontheading{ΣΥΝΤΟΜΟΓΡΑΦΙΕΣ}{Συντομογραφίες}
\renewcommand{\abbrentry}[4]{\item[\textbf{#1}] #3}
\begin{description}[leftmargin=!,labelwidth=1.9cm,labelsep=0.6em,
                    style=sameline,itemsep=0.35em]
\abbrlist
\end{description}


% --- Ενότητα 2η: ΚΥΡΙΩΣ ΚΕΙΜΕΝΟ --------------------------------------
\cleardoublepage
\pagenumbering{arabic}
\setcounter{page}{1}

% =====================================================================
%  ΚΕΦΑΛΑΙΟ 1 — ΕΙΣΑΓΩΓΗ
%  Κάθε αριθμητικό στοιχείο του κεφαλαίου συνοδεύεται από πηγή και από
%  το πειραματικό περιβάλλον στο οποίο παρήχθη. Καμία μέτρηση της
%  παρούσας εργασίας δεν εμφανίζεται σε αυτό το κεφάλαιο.
%
%  Η ΑΓΓΛΙΚΗ ΕΙΝΑΙ ΤΟ ΠΡΩΤΟΤΥΠΟ. Αυτό το αρχείο ακολουθεί το
%  content/en/ch1.tex, μπλοκ προς μπλοκ. Δες docs/SYNC.md.
% =====================================================================

%%blk intro.head from=91d794864574
\chapter{Εισαγωγή: το πρόβλημα της συναίνεσης και το ερώτημα της εργασίας}
\label{ch:intro}

%%blk intro.head.roadmap from=b89a4362725b
Το κεφάλαιο κάνει τέσσερα πράγματα, με αυτή τη σειρά. Διατυπώνει το
πρόβλημα που λύνει η συναίνεση και τα δύο εμπόδια που το κάνουν δύσκολο,
τα σφάλματα και τις καθυστερήσεις. Εξηγεί γιατί δεν το έλυσε οριστικά
ένα πρωτόκολλο, και γιατί οι δημοσιευμένες συγκρίσεις μεταξύ πρωτοκόλλων
δεν μεταφέρονται από τη μία εργασία στην άλλη. Ύστερα δηλώνει τον σκοπό
της εργασίας, τα ερευνητικά ερωτήματα και τα όριά της. Κλείνει με τον
κανόνα που ακολουθεί η εργασία για την ποιότητα των πηγών της, γιατί μια
σύγκριση αξίζει ακριβώς όσο και τα τεκμήρια από κάτω της.

%%blk intro.problem.head from=e3c836bd5a2a
\section{Συμφωνία υπό σφάλματα και καθυστερήσεις}
\label{sec:problem}

%%blk intro.problem.01 from=e1271f592a84
Ένα κατανεμημένο σύστημα αποτελείται από διεργασίες χωρικά
διαχωρισμένες, που επικοινωνούν αποκλειστικά με ανταλλαγή μηνυμάτων. Ένα
σύστημα είναι κατανεμημένο όταν η καθυστέρηση μετάδοσης δεν είναι
αμελητέα σε σχέση με τον χρόνο που μεσολαβεί ανάμεσα σε δύο γεγονότα της
ίδιας διεργασίας \cite{lamport1978time}. Ο ορισμός αυτός έχει μία
συνέπεια: δεν υπάρχει καθολικό ρολόι. Για δύο γεγονότα σε διαφορετικές
διεργασίες είναι ενίοτε αδύνατο να ειπωθεί ποιο συνέβη πρώτο. Η μόνη
εγγενής σχέση διάταξης είναι η μερική σχέση \en{happened-before}
\cite{lamport1978time}. Κάθε υπηρεσία που οφείλει να συμπεριφέρεται σαν
ένα ενιαίο, αξιόπιστο μηχάνημα πρέπει επομένως να κατασκευάσει η ίδια τη
σειρά αυτή, επειδή το δίκτυο δεν της τη δίνει.

%%blk intro.problem.02 from=63ea30011c9d
Η κανονική αφαίρεση με την οποία κατασκευάζεται αυτή η σειρά είναι η
\emph{αναπαραγωγή μηχανής καταστάσεων} (\en{State Machine Replication},
SMR). Η SMR μοντελοποιεί την υπηρεσία ως ντετερμινιστική μηχανή
καταστάσεων και την αναπαράγει σε $n$ κόμβους. Αν όλοι οι σωστοί κόμβοι
εκτελέσουν την ίδια ακολουθία εντολών με την ίδια σειρά, θα καταλήξουν
στην ίδια κατάσταση και θα παραγάγουν τις ίδιες εξόδους
\cite{schneider1990implementing}. Το πρόβλημα της αξιοπιστίας ανάγεται
έτσι σε ένα και μόνο πρόβλημα συντονισμού, τη \emph{συναίνεση}
(\en{consensus}): τη συμφωνία των σωστών κόμβων σε μία κοινή τιμή για
κάθε θέση της ακολουθίας. Το κριτήριο ορθότητας που περιμένει ο πελάτης
από την αναπαραγόμενη υπηρεσία είναι η \emph{γραμμικοποιησιμότητα}
(\en{linearizability}). Κάθε πράξη φαίνεται να συμβαίνει ακαριαία σε
κάποιο σημείο ανάμεσα στην κλήση της (\en{invocation}) και στην απόκρισή
της, και οι πράξεις που δεν επικαλύπτονται χρονικά διατηρούν τη σειρά
τους \cite{herlihy1990linearizability}.

%%blk intro.problem.03 from=3d8d7e1d6b7d
Ένα πρωτόκολλο συναίνεσης κρίνεται από δύο ιδιότητες. Η \emph{ασφάλεια}
(\en{safety}) απαιτεί να μην αποφασίσουν ποτέ δύο σωστοί κόμβοι
διαφορετικά για την ίδια θέση. Η \emph{ζωτικότητα} (\en{liveness})
απαιτεί κάθε αίτημα σωστού πελάτη να αποφασίζεται τελικά. Το θεμελιώδες
εμπόδιο είναι ότι κανένα ντετερμινιστικό πρωτόκολλο δεν εξασφαλίζει τις
δύο ιδιότητες ταυτόχρονα σε πλήρως ασύγχρονο σύστημα. Οι Fischer, Lynch
και Paterson απέδειξαν ότι \emph{κάθε ντετερμινιστικό} πρωτόκολλο για το
πρόβλημα αυτό έχει τη δυνατότητα μη τερματισμού, ακόμη και με μία μόνο
ελαττωματική διεργασία \cite{fischer1985impossibility}. Το αποτέλεσμα είναι δομικό και
όχι τεχνικός περιορισμός: σε ασύγχρονο σύστημα καμία διεργασία δεν μπορεί
να διακρίνει μια νεκρή διεργασία από μια αργή.

%%blk intro.problem.04 from=5f7e2e3ff158
Όλα τα πρωτόκολλα που εξετάζει η εργασία χρησιμοποιούν την ίδια διέξοδο,
το μοντέλο της \emph{μερικής συγχρονίας} (\en{partial synchrony})
\cite{dwork1988consensus}. Στη μορφή που υιοθετούν τα σύγχρονα
πρωτόκολλα, υπάρχει άγνωστη χρονική στιγμή. Είναι η GST, το \en{Global
Stabilization Time}. Μετά την GST, το δίκτυο παραδίδει τα μηνύματα μεταξύ
σωστών κόμβων εντός γνωστού ορίου $\Delta$. Έτσι ακριβώς διατυπώνουν το
μοντέλο τα βυζαντινά πρωτόκολλα της εργασίας \cite{yin2019hotstuff} και
\cite{malkhi2023hotstuff2}\prep. Η ασφάλεια διατηρείται πάντοτε, ακόμη
και πριν την GST, και η ζωτικότητα εξασφαλίζεται μόνο μετά από αυτήν.

%%blk intro.problem.05 from=8a0fcb4b9434
Η συνέπεια για κάθε πειραματική αποτίμηση είναι σημαντική. Ο χρόνος
ανάκαμψης μετά την απώλεια του ηγέτη εξαρτάται σε μεγάλο βαθμό από τη
ρύθμιση των χρονομέτρων, όχι μόνο από τον αλγόριθμο. Δύο εργασίες το
διατυπώνουν ρητά. Η μελέτη αναπαραγωγιμότητας του Raft το σημειώνει
πρώτη \cite{howard2015raftrefloated}. Οι Berger κ.ά.\ το επαναλαμβάνουν
για βυζαντινά συστήματα. Γράφουν ότι ο χρόνος ανάκαμψης εξαρτάται από τις
παραμέτρους των χρονομέτρων και ότι \emph{δεν αποκλείουν} πως στενότερες
ρυθμίσεις θα ήταν εφικτές \cite{berger2023simulation}.

%%blk intro.problem.06 from=1424ac3bc7e6
Η δεύτερη διάσταση που οργανώνει ολόκληρη τη βιβλιογραφία είναι το
\emph{μοντέλο σφαλμάτων}. Στο μοντέλο κατάρρευσης (\en{crash-fault}) ένας
κόμβος μπορεί μόνο να σταματήσει. Στο βυζαντινό μοντέλο ένας κόμβος
μπορεί να συμπεριφερθεί αυθαίρετα, να στείλει αντιφατικά μηνύματα σε
διαφορετικούς παραλήπτες και να συντονιστεί με άλλους ελαττωματικούς
κόμβους \cite{lamport1982byzantine}. Η διάκριση δεν είναι φιλοσοφική.
Καθορίζει πόσοι κόμβοι χρειάζονται, όπως δείχνει η
Ενότητα~\ref{sec:notonce}.

%%blk intro.notonce.head from=cd630a475d44
\section{Γιατί το πρόβλημα δεν λύθηκε μία φορά}
\label{sec:notonce}

%%blk intro.notonce.01 from=e1dc502bef0c
Το Viewstamped Replication δημοσιεύτηκε το 1988. Είναι το πρώτο
πρωτόκολλο της γραμμής που εξετάζει η εργασία. Από τότε έχουν περάσει
σχεδόν τέσσερις δεκαετίες, και η παραγωγή νέων πρωτοκόλλων δεν έχει
σταματήσει. Το Σχήμα~\ref{fig:timeline} δίνει τη χρονολογική σειρά με
έναν ρητό κανόνα: κάθε πρωτόκολλο τοποθετείται στην \emph{πρώτη
δημοσίευσή του με κρίση}, σε πρακτικά συνεδρίου ή σε περιοδικό. Τρία από
τα δέκα πρωτόκολλα δεν έχουν τέτοια δημοσίευση. Το σχήμα τα σχεδιάζει με
κενό δείκτη, τα τοποθετεί στην πρώτη δημόσια διάθεση του σχεδιασμού τους
και σημαίνει την πηγή τους. Η εργασία δηλώνει τον κανόνα εδώ, επειδή η
λέξη «χρονολογικά» χωρίς κανόνα γίνεται αυθαίρετη, και η αυθαιρεσία
φαίνεται.

%%blk intro.notonce.fig.timeline from=057cc358c72b
\dgfigure{fig:timeline}%
{Χρονολογική σειρά των πρωτοκόλλων που εξετάζει η εργασία, με κανόνα
τοποθέτησης την πρώτη δημοσίευση με κρίση.}%
{Χρονολογική σειρά των πρωτοκόλλων που εξετάζει η εργασία, με κανόνα
τοποθέτησης την πρώτη δημοσίευση με κρίση. \textbf{Γεμάτος δείκτης} =
πρώτη δημοσίευση με κρίση, η οποία και ορίζει τη θέση: VR
\cite{oki1988viewstamped}, Paxos \cite{lamport1998parttime}, PBFT
\cite{castro1999pbft}, ZAB \cite{junqueira2011zab}, EPaxos
\cite{moraru2013epaxos}, Raft \cite{ongaro2014raft}, HotStuff
\cite{yin2019hotstuff}. \textbf{Κενός δείκτης} = δεν υπάρχει δημοσίευση
με κρίση. Η θέση δίνεται τότε από την πρώτη δημόσια διάθεση του
σχεδιασμού, και το σχήμα τη σημαίνει. KRaft: KIP-500, τεχνική
τεκμηρίωση\grey\ του 2019 \cite{kafka2019kip500}. Η γενική διάθεση ήρθε
αργότερα, με το Kafka~3.3· η χρονολογία εκείνης της έκδοσης δεν
τεκμηριώνεται από την πηγή που παρατίθεται εδώ, και δεν ορίζει τη θέση
στο σχήμα. Tendermint: αλγοριθμικό \en{preprint}\prep\ του
2018 \cite{buchman2018gossip}, \en{whitepaper}\grey\ του 2014
\cite{kwon2014tendermint}. Κριθείσα εργασία υπάρχει
\cite{cason2021tendermint}, αλλά εμφανίστηκε το 2021 και περιγράφει τον
σχεδιασμό, την αρχιτεκτονική και τις επιδόσεις του \emph{δικτύου}·
τοποθετώντας εκεί το πρωτόκολλο θα το βάζαμε επτά χρόνια μετά τον
σχεδιασμό που υλοποιεί. Το σχήμα τοποθετεί επομένως το Tendermint στο
αλγοριθμικό του \en{preprint} με κενό δείκτη, και η παράγραφος αυτή
είναι η επιφύλαξη που συνοδεύει την επιλογή. HotStuff-2: εκτεταμένη
περίληψη σε ePrint\prep\ \cite{malkhi2023hotstuff2}. \emph{Δύο
επιφυλάξεις προέλευσης.} Το \cite{lamport1998parttime} φέρει ημερομηνία
παραλαβής τον Ιανουάριο του 1990 και αποδοχής τον Μάρτιο του 1998, αλλά
τα δύο αντίτυπα του \en{corpus} τη \emph{διατυπώνουν} διαφορετικά (βλ.\
\S\ref{sec:sources}). Η τεχνική αναφορά SRC~RR-49 χρονολογείται ευρέως
στο 1989. Δεν την εντοπίσαμε ως τεκμήριο και γι' αυτό δεν ορίζει τη θέση.
Το PBFT έχει και εκτενή εκδοχή σε περιοδικό το 2002
\cite{castro2002pbftrecovery}. Θεωρητικά ορόσημα:
\cite{lamport1978time,lamport1982byzantine,fischer1985impossibility,%
dwork1988consensus,schneider1990implementing,herlihy1990linearizability}.}%
{% fig-timeline — first-refereed-publication rule, CFT above the axis, BFT below.
% Structure figure: dgflow. The only geometry it sets is the year-to-length
% mapping, and that is tied to \dgwidth so the timeline is exactly as wide as
% every other figure in the thesis: 51 years (1977-2028) across \dgwidth.
% Every human-readable string is a label key; protocol names, venues and years
% are identical in both editions and stay literal.
\begin{dgflow}[x={\dgwidth/51}]
  \def\yzero{1977}

  % Theoretical milestones. A third category — neither CFT nor BFT protocols —
  % so dgAux rather than either series colour.
  \node[dgnote, text=dgAux, anchor=west] at (0,4.45)
    {\textbf{\fl{timeline.milestones.title}}\\
     1978 \fl{timeline.milestone.clocks}\, ·\, 1982 \fl{timeline.milestone.byzmodel}\, ·\, 1985 FLP\\
     1988 \fl{timeline.milestone.partialsync}\, ·\, 1990 SMR\, ·\, 1990 \fl{timeline.milestone.linearizability}};

  % The time axis and its five-year ticks. The year LABELS are drawn last, so
  % that their dgtag background sits over the leader lines rather than under.
  \draw[dgaxis] (0,0) -- (51,0);
  \foreach \yr in {1980,1985,1990,1995,2000,2005,2010,2015,2020,2025}{
    \pgfmathsetmacro{\xx}{\yr-\yzero}
    \draw[dgaxis,-] (\xx,-0.08) -- (\xx,0.08);}

  % Which half is which. At the left end, where the axis is still empty.
  \node[dgtag, fill=dgPaper, text=dgCFT, anchor=west] at (0.4, 1.00) {\textbf{CFT}};
  \node[dgtag, fill=dgPaper, text=dgBFT, anchor=west] at (0.4,-1.00) {\textbf{BFT}};

  % Crash-fault protocols, above the axis. Three stagger levels: 2011-2019 is
  % 1.7 cm of axis carrying 5.6 cm of label, and the anchors are chosen so that
  % no leader line ever runs through another protocol's label.
  % year / level / anchor / filled marker? / label
  \foreach \yr/\lvl/\anc/\solid/\lbl in {%
      1988/1.00/{south}/{1}/{VR\\PODC '88},%
      1998/1.00/{south}/{1}/{Paxos\\TOCS '98},%
      2011/1.00/{south east}/{1}/{ZAB\\DSN '11},%
      2013/1.95/{south east}/{1}/{EPaxos\\SOSP '13},%
      2014/2.90/{south}/{1}/{Raft\\ATC '14},%
      2019/1.00/{south west}/{0}/{KRaft\\KIP-500 '19\grey}}{
    \pgfmathsetmacro{\xx}{\yr-\yzero}
    \draw[dgarrow,-,draw=dgCFT] (\xx,0.07) -- (\xx,\lvl);
    \ifnum\solid=1 \fill[dgCFT] (\xx,0.07) circle (1.7pt);
    \else \draw[dgarrow,-,draw=dgCFT,fill=dgPaper] (\xx,0.07) circle (1.7pt);\fi
    \node[dgtag, fill=dgPaper, text=dgCFT, anchor=\anc] at (\xx,\lvl) {\lbl};}

  % Byzantine-fault protocols, below the axis. Same three levels, same reason:
  % 2018, 2019 and 2023 are 1.1 cm apart and their labels are 1.5 cm wide.
  \foreach \yr/\lvl/\anc/\solid/\lbl in {%
      1999/-1.00/{north}/{1}/{PBFT\\OSDI '99},%
      2018/-1.00/{north east}/{0}/{Tendermint\\arXiv '18\prep},%
      2019/-1.95/{north east}/{1}/{HotStuff\\PODC '19},%
      2023/-2.90/{north}/{0}/{HotStuff-2\\ePrint '23\prep}}{
    \pgfmathsetmacro{\xx}{\yr-\yzero}
    \draw[dgarrow,-,draw=dgBFT] (\xx,-0.07) -- (\xx,\lvl);
    \ifnum\solid=1 \fill[dgBFT] (\xx,-0.07) circle (1.7pt);
    \else \draw[dgarrow,-,draw=dgBFT,fill=dgPaper] (\xx,-0.07) circle (1.7pt);\fi
    \node[dgtag, fill=dgPaper, text=dgBFT, anchor=\anc] at (\xx,\lvl) {\lbl};}

  % Year labels last: 1988, 1999, 2011, 2014 and 2019 all sit within half a
  % label of a five-year tick, so something has to be on top and it should be
  % the year.
  \foreach \yr in {1980,1985,1990,1995,2000,2005,2010,2015,2020,2025}{
    \pgfmathsetmacro{\xx}{\yr-\yzero}
    \node[dgtag, fill=dgPaper, text=dgGrey, anchor=north] at (\xx,-0.16) {\yr};}
\end{dgflow}
}

%%blk intro.notonce.02 from=d35f19dfee55
Η επανάληψη δεν οφείλεται σε αδυναμία των προηγούμενων λύσεων να είναι
ορθές. Οφείλεται στο ότι η ορθότητα είναι μόνο η μία διάσταση. Κάθε
πρωτόκολλο είναι μια \emph{θέση} σε έναν χώρο σχεδίασης, και κάθε
βελτίωση εκεί αγοράζεται με συγκεκριμένο τίμημα. Η πρόσφατη ενοποιητική
εργασία των Amiri κ.ά.\ το διατυπώνει ρητά. Περιγράφει τα βυζαντινά
πρωτόκολλα ως σημεία ενός χώρου με περισσότερες από δέκα διαστάσεις
σχεδίασης. Στον χώρο αυτόν κάθε επιλογή είναι μια ένα-προς-ένα
συνάρτηση από πρωτόκολλο σε πρωτόκολλο \cite{amiri2024bedrock}.

%%blk intro.notonce.03 from=11de5b63b6e0
Η σαφέστερη ένδειξη αυτού είναι η αριθμητική των \en{quorums}
(Σχήμα~\ref{fig:quorum}). Το πλήθος των κόμβων δεν αυξάνεται αυθαίρετα.
Μετακινείται μόνο όταν αγοράζει κάτι συγκεκριμένο και κατονομάσιμο. Με
σφάλματα κατάρρευσης αρκεί ένα πλειοψηφικό \en{quorum}. Μια συστάδα
$n = 2f+1$ κόμβων παραμένει διαθέσιμη όσο λειτουργεί μια πλειοψηφία, άρα
ανέχεται $f$ σφάλματα \cite{ongaro2014raft}. Η γενική διατύπωση της
ασφάλειας ως τομή πλειοψηφικών συνόλων ανάγεται στο
\cite{lamport1998parttime}.

%%blk intro.notonce.04 from=cb4943a00b1a
Με βυζαντινά σφάλματα και μη αυθεντικοποιημένα μηνύματα απαιτούνται
τουλάχιστον $n = 3f+1$ κόμβοι. Το φράγμα διατυπώνεται πρώτη φορά το 1980,
ως πρόβλημα διαδραστικής συνέπειας \cite{pease1980reaching}. Η αλληγορία
των βυζαντινών στρατηγών είναι μεταγενέστερη \cite{lamport1982byzantine}.
Το ίδιο φράγμα ισχύει και για τα πρωτόκολλα μερικής συγχρονίας που
εξετάζονται εδώ \cite{dwork1988consensus,amiri2024bedrock}. Εκεί, όμως,
δεν προκύπτει από την απουσία υπογραφών, και η
Ενότητα~\ref{sec:bounds} δείχνει γιατί.

%%blk intro.notonce.05 from=c7c6ac12e911
Το έμπιστο υλικό (\en{trusted execution environment}) εμποδίζει έναν
ελαττωματικό κόμβο να στείλει αντιφατικά μηνύματα χωρίς να γίνει
αντιληπτός. Αν προστεθεί τέτοιο υλικό, το φράγμα επιστρέφει στο $2f+1$
\cite{amiri2024bedrock}. Η εξοικονόμηση μιας φάσης επικοινωνίας κόστιζε
αρχικά $5f+1$ κόμβους \cite{martin2006fab}\unver. Με αυθεντικοποίηση αρκούν
$5f-1$. Ο αριθμός αυτός είναι και το αποδεδειγμένο κάτω φράγμα για
βυζαντινή συναίνεση δύο βημάτων \cite{kuznetsov2021optimal}\unver, και
επιτυγχάνεται στην πράξη \cite{amiri2024bedrock}. Τέλος, η υποστήριξη
δίκαιης διάταξης (\en{order-fairness}) απαιτεί τουλάχιστον $4f+1$ κόμβους
στη \emph{λιγότερο} απαιτητική παραμετροποίησή της. Καθώς η παράμετρος
$\gamma$ πλησιάζει το $0{,}5$ και η απαίτηση δικαιοσύνης γίνεται
αυστηρότερη, το φράγμα $n > 4f/(2\gamma-1)$ αυξάνεται απεριόριστα
\cite{kelkar2020orderfairness,amiri2024bedrock}\unver.

%%blk intro.notonce.fig.quorum from=5307770dd08c
\dgfigure{fig:quorum}%
{Αριθμητική \protect\en{quorum}: το απαιτούμενο πλήθος κόμβων $n$ ως
συνάρτηση των ανεκτών σφαλμάτων $f$.}%
{Αριθμητική \protect\en{quorum}: το απαιτούμενο πλήθος κόμβων $n$ ως
συνάρτηση των ανεκτών σφαλμάτων $f$. Δεν πρόκειται για μετρήσεις. Είναι
καθαρή αριθμητική από φράγματα που δηλώνονται στη βιβλιογραφία. Για
$f=3$: η βυζαντινή ανοχή κοστίζει 10 κόμβους έναντι 7, η δίκαιη διάταξη
13 και η εξοικονόμηση μίας φάσης 16. \emph{Πηγές.} Το
\cite{amiri2024bedrock} συγκεντρώνει και τα πέντε φράγματα: $3f{+}1$,
$2f{+}1$ και $5f{\pm}1$ στο \S3.2 E1, έμπιστο υλικό στο Παράρτ.~A, O7,
και δίκαιη διάταξη στο \S3.3 Q1 και στη σχεδιαστική επιλογή 13. Για τα
φράγματα αυτά είναι δευτερογενής πηγή. Για το $3f{+}1$ δίνει πέντε
αναφορές, από τις οποίες η εργασία κατονομάζει μία, το
\cite{lamport1982byzantine}· οι άλλες τέσσερις είναι Bracha 1984,
Bracha και Toueg 1985, Correia κ.ά.\ 2006 και \cite{dwork1988consensus}.
Η πρωτογενής διατύπωση είναι ακόμη παλαιότερη, του 1980
\cite{pease1980reaching}, το οποίο το \cite{amiri2024bedrock} \emph{δεν}
παραθέτει --- η απόδοση αυτή είναι της ίδιας της εργασίας, και το
σκεπτικό της δίνεται στην Ενότητα~\ref{sec:bounds}. Για τα υπόλοιπα
φράγματα η ανασκόπηση αποδίδει το $5f{+}1$ στο
\cite{martin2006fab}\unver, το $5f{-}1$ στα \cite{kuznetsov2021optimal}\unver\
και \cite{abraham2021goodcase}\unver\ (παραθέτει και τα δύο), και το
$4f{+}1$ στο \cite{kelkar2020orderfairness}\unver.
Από αυτές τις πηγές, μόνο τα \cite{pease1980reaching} και
\cite{lamport1982byzantine} επαληθεύτηκαν ως προς τα βιβλιογραφικά τους
στοιχεία, μέσω ευρετηρίου. Τα υπόλοιπα τέσσερα δεν επαληθεύτηκαν και τα
στοιχεία τους προέρχονται από τον κατάλογο αναφορών του
\cite{amiri2024bedrock}. \emph{Κανένα από τα έξι δεν βρίσκεται στο
\en{corpus}.} \emph{Σημείωση ακρίβειας.} Το \cite{amiri2024bedrock}
διακρίνει το μέγεθος του \en{quorum} από το μέγεθος του δικτύου. Σε
σχεδιάσεις με ενεργούς και παθητικούς κόμβους το \en{quorum} πέφτει στο
$2f{+}1$ ενώ το $n$ παραμένει $3f{+}1$.}%
{% fig-quorum — required n as a function of tolerated faults f. Pure arithmetic
% from bounds stated in the literature; no measurements.
% Data figure: dgdata. The five series are plotted in the order that makes the
% standard cycle list hand them their meanings: dgCFT crash, dgBFT Byzantine,
% then dgAux / dgSand / dgPurple for the three further bounds.
\begin{dgdata}[
  xlabel={\fl{quorum.xlabel} $f$},
  ylabel={\fl{quorum.ylabel} $n$},
  xmin=0.7, xmax=5.3, ymin=0, ymax=28,
  xtick={1,2,3,4,5}, ytick={0,5,10,15,20,25},
  legend style={at={(0.5,-0.30)}, anchor=north}, legend columns=1,
]
\addplot coordinates {(1,3)(2,5)(3,7)(4,9)(5,11)};
  \addlegendentry{$2f{+}1$ — \fl{quorum.crash}}
\addplot coordinates {(1,4)(2,7)(3,10)(4,13)(5,16)};
  \addlegendentry{$3f{+}1$ — \fl{quorum.byzantine}}
\addplot coordinates {(1,5)(2,9)(3,13)(4,17)(5,21)};
  \addlegendentry{$4f{+}1$ — \fl{quorum.fairorder} ($\gamma=1$)}
\addplot coordinates {(1,6)(2,11)(3,16)(4,21)(5,26)};
  \addlegendentry{$5f{+}1$ — \fl{quorum.onephaseless} (FaB)}
\addplot coordinates {(1,4)(2,9)(3,14)(4,19)(5,24)};
  \addlegendentry{$5f{-}1$ — \fl{quorum.twostep}}
\end{dgdata}
}

%%blk intro.notonce.06 from=c2336f35c2a2
Η ίδια λογική διέπει τη δεύτερη μεγάλη διάσταση, την \emph{αλλαγή όψης}
(\en{view change}). Στη διαδρομή χωρίς σφάλματα τα πρωτόκολλα διαφέρουν
κυρίως ως προς την τοπολογία και το πλήθος των φάσεων. Αποκλίνουν
δραστικότερα στην αντικατάσταση του ηγέτη, που οι σχεδιαστές του HotStuff
ονομάζουν «\en{the epicenter of replication}» \cite{yin2019hotstuff}. Ο
Πίνακας~\ref{tab:viewchange} συνοψίζει τέσσερις χαρακτηριστικές θέσεις σε
αυτόν τον άξονα.

%%blk intro.notonce.tab.viewchange from=b62f864471a2
\begin{table}[t]
\centering
\caption[Τέσσερις χαρακτηριστικές θέσεις στον άξονα της αλλαγής όψης:
κόστος, αποκρισιμότητα και πλήθος φάσεων.]%
{Τέσσερις χαρακτηριστικές θέσεις στον άξονα της αλλαγής όψης:
κόστος, αποκρισιμότητα και πλήθος φάσεων. Ο πίνακας αναπαράγει τον
Πίνακα~1 του \cite{malkhi2023hotstuff2}\prep\ μαζί με την επιφύλαξη που
η ίδια πηγή διατυπώνει. Δεν πρέπει να διαβαστεί χωρίς αυτήν. Δεν είναι
εξαντλητικός, επειδή η ίδια πηγή απαριθμεί άλλες σχεδιάσεις στον ίδιο
άξονα.}
\label{tab:viewchange}
\small
\begin{tabular}{@{}l c c c@{}}
\toprule
\textbf{Πρωτόκολλο} & \textbf{Κόστος αλλαγής όψης} &
\textbf{Αποκρίσιμη (\en{responsive})} & \textbf{Φάσεις} \\
\midrule
PBFT \cite{castro1999pbft}        & $O(n^{2})$ & ναι  & 2 \\
Tendermint \cite{buchman2018gossip}\prep & $O(n)$   & όχι  & 2 \\
HotStuff \cite{yin2019hotstuff}   & $O(n)$     & ναι  & 3 \\
HotStuff-2 \cite{malkhi2023hotstuff2}\prep & $O(n)$ & ναι & 2 \\
\bottomrule
\end{tabular}

\vspace{4pt}
\begin{minipage}{0.94\linewidth}\setstretch{1.08}\footnotesize
\textbf{Επιφύλαξη της ίδιας της πηγής.} Η πηγή
\cite{malkhi2023hotstuff2}\prep\ σημειώνει ρητά ότι τα PBFT και
Tendermint έχουν κόστος αλλαγής όψης $O(n^{3})$ και $O(n^{2})$
αντίστοιχα, και ότι τα φράγματα του πίνακα προκύπτουν αφού και τα δύο
βελτιωθούν κατά παράγοντα $O(n)$ --- βελτίωση που η ίδια πηγή αποδίδει
αλλού στη συνάθροιση υπογραφών. Η αναλυτική μέτρηση πολυπλοκότητας στο
\cite{yin2019hotstuff} (Πίνακας~1) δίνει για το PBFT $O(n^{2})$
αυθεντικοποιητές με σταθερό ηγέτη, $O(n^{3})$ σε αλλαγή όψης και
$O(f n^{3})$ σε αλληλουχία $f$ αποτυχιών ηγέτη. Δεν πρόκειται για
ανεξάρτητη επιβεβαίωση, καθώς οι δύο εργασίες μοιράζονται συγγραφέα. Το
\cite{jalalzai2023fasthotstuff} επιστρέφει σε δύο γύρους διατηρώντας την
αποκρισιμότητα, αλλά δηλώνει ρητά ότι \emph{δεν} αποφεύγει τη μεταφορά
τετραγωνικού πλήθους μηνυμάτων αλλαγής όψης όταν ο ηγέτης αποτύχει. Αυτό
που μειώνει είναι το πλήθος των υπογραφών που επαληθεύει κάθε κόμβος, από
$O(n)$ σε δύο. Οι Malkhi και Nayak δημοσιεύουν εκτεταμένη περίληψη σε
ePrint, που δεν έχει περάσει από κρίση \cite{malkhi2023hotstuff2}\prep.
Αναφέρεται ως τέτοια.

\smallskip
\textbf{Φάση δεν σημαίνει καθυστέρηση μηνύματος.} Ο πίνακας μετρά
\emph{φάσεις}, με τη σύμβαση της πηγής του. Μια φάση του HotStuff
κοστίζει δύο διαδοχικές καθυστερήσεις, επειδή περνά από τον ηγέτη. Μια
φάση ψηφοφορίας του PBFT κοστίζει μία, επειδή είναι όλοι-προς-όλους. Οι
δύο μονάδες δεν συγκρίνονται, και η διάκριση ορίζεται μία φορά στην
Ενότητα~\ref{sec:performance}.
\end{minipage}
\end{table}

%%blk intro.notonce.07 from=21b30534c3c2
Ο πίνακας δείχνει ότι κανένα από τα τέσσερα πρωτόκολλα δεν είναι
«καλύτερο». Το καθένα πληρώνει διαφορετικό λογαριασμό. Το Tendermint
καταργεί το υποπρωτόκολλο αλλαγής όψης και πληρώνει με αναμονή του
χειρότερου δυνατού $\Delta$. Δηλαδή πληρώνει με την \emph{αποκρισιμότητα}
(\en{responsiveness}), την ιδιότητα να προχωρά το πρωτόκολλο με τον
πραγματικό και όχι με τον μέγιστο χρόνο του δικτύου
\cite{buchman2018gossip}\prep\ και \cite{malkhi2023hotstuff2}\prep. Το
HotStuff πληρώνει με τρίτη φάση \cite{yin2019hotstuff}. Το HotStuff-2
ισχυρίζεται ότι κανένα από τα δύο τιμήματα δεν ήταν αναγκαίο
\cite{malkhi2023hotstuff2}\prep.

%%blk intro.opengap.head from=9649f2a47d7e
\section{Γιατί η σύγκριση παραμένει ανοιχτό ερώτημα}
\label{sec:opengap}

%%blk intro.opengap.01 from=9d0602b13d5f
Ο χώρος σχεδίασης είναι γνωστός και οι πολυπλοκότητες υπολογισμένες. Θα
μπορούσε επομένως κανείς να υποθέσει ότι η σύγκριση των πρωτοκόλλων είναι
λυμένο πρόβλημα. Δεν είναι, για τρεις λόγους που τεκμηριώνονται από την
ίδια τη βιβλιογραφία.

%%blk intro.opengap.02 from=2495649c056e
\paragraph{Πρώτον: τα δημοσιευμένα νούμερα δεν μεταφέρονται.}
Το Σχήμα~\ref{fig:scaling} το δείχνει στη δυσκολότερη δυνατή περίπτωση
για την ένσταση: ίδια εργασία, ίδια ενιαία υλοποίηση, ίδια μετρική, ίδιο
εύρος $n$. Αλλάζει μόνο η τοπολογία ανάπτυξης. Το κόστος κλιμάκωσης του
PBFT περνά από 65\% σε 84\%, του PoE από 39\% σε 51\% και του FTB από
32\% σε 36\% \cite{amiri2024bedrock}. Η ίδια εργασία δείχνει και
αντιστροφή κατάταξης χωρίς καμία αλλαγή περιβάλλοντος. Το SBFT αποδίδει
12\% \emph{υψηλότερο} \en{throughput} από το HotStuff σε $n=4$ και 8\%
\emph{χαμηλότερο} σε $n=100$, στο ίδιο πείραμα \cite{amiri2024bedrock}.

%%blk intro.opengap.03 from=780afdd88510
Όταν αλλάξει και η υλοποίηση, οι διαφορές γίνονται τάξεις μεγέθους. Σε
προσομοίωση πλανητικού δικτύου με $n=256$ κόμβους, το Kauri αποδίδει
περίπου εικοσαπλάσιο \en{throughput} από το HotStuff-SECP256K1, δηλαδή
4.917 πράξεις ανά δευτερόλεπτο έναντι 246. Εδώ κάθε πρωτόκολλο έχει
διαφορετική υλοποίηση, και οι τιμές είναι προσομοιωμένες και όχι
μετρημένες \cite{berger2023simulation}.

%%blk intro.opengap.fig.scaling from=bca3f3287031
\dgfigure{fig:scaling}%
{Το ίδιο πείραμα σε δύο τοπολογίες ανάπτυξης: ποσοστιαία μείωση του
\protect\en{throughput} καθώς οι κόμβοι αυξάνονται από 4 σε 100.}%
{Το ίδιο πείραμα σε δύο τοπολογίες ανάπτυξης: ποσοστιαία μείωση του
\protect\en{throughput} καθώς οι κόμβοι αυξάνονται από 4 σε 100. Όλα τα
νούμερα προέρχονται από την ίδια εργασία και την ίδια ενιαία υλοποίηση
των πρωτοκόλλων \cite{amiri2024bedrock}. \textbf{(α)}~Ανάπτυξη σε ένα
\en{datacenter}: Amazon EC2, στιγμιότυπα \texttt{c4.2xlarge} (8 vCPU,
15\,GB RAM, Intel Xeon E5-2666~v3 στα 3{,}50\,GHz), ένας κόμβος ανά
εικονική μηχανή, μέγεθος παρτίδας 400, ωφέλιμο φορτίο 128/128\,\en{byte},
πελάτες σε κλειστό βρόχο \cite{schroeder2006openclosed}\unver, μέσος όρος πέντε
εκτελέσεων, χωρίς σφάλματα. Τρεις διευκρινίσεις για το σχήμα. Το FaB
απαιτεί $5f{+}1$, άρα ξεκινά από $n=6$. Το \cite{amiri2024bedrock}
υλοποιεί το HotStuff αλυσιδωτά, με \en{pipelining}. Τα FTB και FLB είναι
νέα πρωτόκολλα που εισάγει η ίδια εργασία, όχι καθιερωμένα.
\textbf{(β)}~Τα τρία πρωτόκολλα για τα οποία η ίδια εργασία δίνει και
γεωκατανεμημένο νούμερο (Παράρτ.~D, ίδιο πλαίσιο υλοποίησης, ίδια
μετρική, ίδια παρτίδα 400, χωρίς σφάλματα). \emph{Δύο επιφυλάξεις που
οφείλει να δηλώσει η ίδια η εργασία.} Στο γεωκατανεμημένο σύνολο το
\cite{amiri2024bedrock} αυξάνει το \en{pipelining stretch} των Kauri και
FTB από 4 σε 6, και το εύρος του FTB φτάνει στο $n=99$ αντί για $n=100$.
Καμία από τις δύο δεν αφορά το PoE, άρα οι καθαρές --- χωρίς καμία άλλη
μεταβολή --- συγκρίσεις είναι δύο: του \textbf{PBFT} (65\,\% $\to$
84\,\%) και του \textbf{PoE} (39\,\% $\to$ 51\,\%). Και για τα τρία,
πάντως, ο αλγόριθμος και η μετρική παραμένουν ίδιοι, και αυτό που αλλάζει
είναι η τοπολογία ανάπτυξης.}%
{% fig-scaling — the same experiment, two deployment topologies. Two panels, so
% two dgbarchart environments stacked: the standard fixes every figure at
% \dgwidth, which leaves no room to set them side by side.
% All seven protocols are Byzantine, so the LAN series is dgBFT in both panels.
% Panel (b)'s second series is a second DEPLOYMENT of the same Byzantine
% protocols, so it takes a further-series colour (dgAux) — dgCFT would claim
% these are crash-fault protocols.
\begin{dgbarchart}[
  title={\fl{scaling.panel.lan}},
  ylabel={\fl{scaling.ylabel} (\%)},
  ymax=78,
  symbolic x coords={PBFT,FaB,SBFT,HotStuff,PoE,FTB,Kauri}, xtick=data,
  nodes near coords, every node near coord/.append style={dgtag, fill=dgPaper},
]
\addplot[dgBFT, fill=dgBFT] coordinates
  {(PBFT,65) (FaB,63) (SBFT,55) (HotStuff,45) (PoE,39) (FTB,32) (Kauri,31)};
\end{dgbarchart}

\par\vspace{5mm}

\begin{dgbarchart}[
  title={\fl{scaling.panel.geo}},
  ylabel={\fl{scaling.ylabel} (\%)},
  ymax=104,
  symbolic x coords={PBFT,PoE,FTB}, xtick=data,
  legend pos=north east,
  nodes near coords, every node near coord/.append style={dgtag, fill=dgPaper},
]
\addplot[dgBFT, fill=dgBFT,
         every node near coord/.append style={xshift=-2pt}]
  coordinates {(PBFT,65) (PoE,39) (FTB,32)};
  \addlegendentry{LAN}
\addplot[dgAux, fill=dgAux,
         every node near coord/.append style={xshift=2pt}]
  coordinates {(PBFT,84) (PoE,51) (FTB,36)};
  \addlegendentry{\fl{scaling.legend.geo}}
\end{dgbarchart}
}

%%blk intro.opengap.04 from=9a4a2e527cee
\paragraph{Δεύτερον: η ίδια σύγκριση αντιστρέφεται ανάλογα με το τι
μετριέται.}
Το πιο διδακτικό παράδειγμα είναι το EPaxos. Η πρωτότυπη εργασία
ισχυρίζεται βέλτιστη καθυστέρηση δέσμευσης σε δίκτυο ευρείας περιοχής,
ομοιόμορφη κατανομή φορτίου και ομαλή υποβάθμιση, με πειράματα σε Amazon
EC2 \cite{moraru2013epaxos}. Η ανεξάρτητη επαναξιολόγηση των Tollman
κ.ά.\ καταλήγει σε τρία συμπεράσματα \cite{tollman2021epaxosrevisited}.
\emph{Επιβεβαιώνει} τον ισχυρισμό για τη \emph{διάμεση} καθυστέρηση.
Δείχνει όμως καθυστέρηση ουράς (\en{tail latency}) πάνω από τέσσερις
φορές χειρότερη από το Multi-Paxos. Δείχνει επίσης ότι το ποσοστό των
πράξεων που ολοκληρώνονται σε έναν κύκλο μετ' επιστροφής έπεσε στο 65\%
σε ένα από τα πειράματά τους, υπό φόρτο με υψηλή σύγκρουση εντολών τον
οποίο η πρωτότυπη μεθοδολογία δεν προκαλούσε. Οι Matte κ.ά.\ αλλάζουν
πάλι τη μετρική, από απόλυτο \en{throughput} σε αποδοτικότητα πόρων.
Καταλήγουν ότι σε περιβάλλον σταθερού προϋπολογισμού πόρων, τυπικό των
πραγματικών νεφών, το μη βελτιστοποιημένο Multi-Paxos επιτυγχάνει
\emph{σχεδόν} διπλάσιο \en{throughput} από το EPaxos
\cite{matte2021wasteful}. Και οι τρεις εργασίες είναι σωστές. Απλώς δεν
μετρούν το ίδιο πράγμα. Η συνέπεια για κάθε συγκριτική αποτίμηση είναι
ότι η επιλογή της μετρικής δεν είναι τεχνική λεπτομέρεια αλλά μέρος του
συμπεράσματος, και πρέπει να δηλώνεται πριν από τη μέτρηση.

%%blk intro.opengap.05 from=84898cb359c8
\paragraph{Τρίτον: η διαφορά του πρωτοκόλλου δεν βρίσκεται εκεί που την
ψάχνει η μέτρηση.}
Οι Howard και Mortier αναλύουν τα δύο πρωτόκολλα που κυριαρχούν στα
παραγωγικά συστήματα. Καταλήγουν ότι Paxos και Raft ακολουθούν «πολύ
παρόμοια προσέγγιση, διαφέροντας μόνο ως προς την εκλογή ηγέτη»
\cite{howard2020paxosraft}. Αν αυτό ισχύει, τότε μια μέτρηση που
συγκρίνει δύο \emph{συστήματα} τα οποία υλοποιούν αυτά τα πρωτόκολλα
μετράει, στη διαδρομή χωρίς σφάλματα, μηχανική υλοποίησης και όχι
αλγοριθμική διαφορά. Μια τέτοια μέτρηση επομένως \emph{απομονώνει} την
υλοποίηση, που είναι ακριβώς το σωστό πείραμα αν αυτό είναι το ερώτημα.
Στη διαδρομή αλλαγής ηγέτη ισχύει το αντίστροφο. Εκεί συγκεντρώνεται όλη
η αλγοριθμική διαφορά, και οι ίδιοι οι συγγραφείς κάνουν δύο χωριστές
δηλώσεις. Η πρώτη είναι δομική διαπίστωση: η φάση εκλογής του Raft είναι
\emph{ελαφρύτερη σε μηνύματα}, επειδή δεν μεταφέρει εγγραφές του
\en{log}. Η δεύτερη είναι ρητή πρόβλεψη: η εκλογή αναμένεται \emph{και
πιο αργή και με μεγαλύτερη διακύμανση} στον χρόνο
\cite{howard2020paxosraft}. Δύο επιφυλάξεις είναι υποχρεωτικές. Η
σύγκριση γίνεται πάνω σε μια \emph{απλοποιημένη} διατύπωση του Paxos στην
ορολογία του Raft, και οι συγγραφείς δηλώνουν ότι και οι δύο αλγόριθμοι
που παρουσιάζουν είναι αφελείς εκ σχεδιασμού \cite{howard2020paxosraft}.

%%blk intro.opengap.06 from=42c7e04b3403
Το κενό, τέλος, έχει διατυπωθεί ρητά, και έχει καλυφθεί μερικώς. Οι
Ailijiang κ.ά.\ διαπιστώνουν ότι τα πρωτόκολλα ισχυρής συνέπειας
«εμφανίζουν ευρέως μεταβαλλόμενες επιδόσεις υπό διαφορετικές συνθήκες
δικτύου, φόρτου και ανάπτυξης, και καμία προηγούμενη μελέτη δεν προσφέρει
ολοκληρωμένη ανατομία και σύγκριση των επιδόσεών τους» (μετάφραση δική
μας), και το καλύπτουν για την οικογένεια κατάρρευσης
\cite{ailijiang2019dissecting}. Οι Amiri κ.ά.\ κάνουν το αντίστοιχο για
τη βυζαντινή \cite{amiri2024bedrock}. Οι Alqahtani και Demirbas μετρούν
πρωτόκολλα και των δύο οικογενειών στο ίδιο πλαίσιο
\cite{alqahtani2021bottlenecks}.

%%blk intro.opengap.07 from=4c5e67029ab6
Αυτό που παραμένει ανοιχτό δεν είναι επομένως η ενιαία μέτρηση καθαυτή.
Είναι η μέτρηση των \emph{παραγωγικών υλοποιήσεων} --- αυτών που
πραγματικά αναπτύσσονται --- με ενιαίο όργανο, ενιαίο μοντέλο φόρτου και
σενάρια σφάλματος. Και οι τρεις παραπάνω εργασίες μετρούν ερευνητικά
πρωτότυπα, ή ενιαίες επανυλοποιήσεις μέσα σε ένα πλαίσιο, και όχι τα
συστήματα που τρέχουν σε παραγωγή. Όπως δείχνει το
Σχήμα~\ref{fig:scaling}, η επανυλοποίηση είναι ακριβώς η μεταβλητή που
μπορεί να μετακινήσει το αποτέλεσμα.

%%blk intro.aims.head from=ce135bb1be78
\section{Σκοπός, ερευνητικά ερωτήματα και όρια της εργασίας}
\label{sec:aims}

%%blk intro.aims.01 from=1d3ab27b5434
Σύμφωνα με την εγκεκριμένη πρόταση εκπόνησης (βλ.\ Παράρτημα~Α), η
εργασία έχει τρεις σκοπούς. Ο πρώτος είναι η μελέτη της εξέλιξης των
πρωτοκόλλων συναίνεσης. Ο δεύτερος είναι η θεωρητική και αρχιτεκτονική
ανάλυση των κυριότερων από αυτά. Ο τρίτος είναι η σύγκριση των επιδόσεων
επιλεγμένων υλοποιήσεων σε αντιπροσωπευτικά σενάρια. Τα κριτήρια της
σύγκρισης είναι τέσσερα: η καθυστέρηση συναίνεσης, το \en{throughput}, η
ανοχή σε σφάλματα και η κλιμάκωση. Τα ερευνητικά ερωτήματα που
ακολουθούν εξειδικεύουν τους σκοπούς αυτούς.

%%blk intro.aims.rq from=5fdf024d64a8
\begin{description}[leftmargin=2.1em,style=nextline,font=\normalfont\bfseries]
  \item[ΕΕ1 --- Εξέλιξη.] Ποιες συγκεκριμένες ανταλλαγές (\en{trade-offs})
  εξηγούν κάθε μετακίνηση μέσα στον χώρο σχεδίασης, και ποιο τίμημα
  πληρώνεται σε κάθε μία; Απαντάται αναλυτικά, από τις πρωτογενείς πηγές.
  \item[ΕΕ2 --- Κοινή δομή.] Ποια αναλλοίωτα (\en{invariants}) μοιράζονται
  όλα τα πρωτόκολλα, ώστε να μπορούν να περιγραφούν σε ενιαίο σχήμα αντί
  για κατάλογο ανεξάρτητων αλγορίθμων;
  \item[ΕΕ3 --- Συγκρισιμότητα.] Υπό ποιες προϋποθέσεις είναι θεμιτή η
  απευθείας σύγκριση δύο υλοποιήσεων, και πότε η μέτρηση αποτυπώνει
  μηχανική υλοποίησης αντί για ιδιότητα του πρωτοκόλλου;
  \item[ΕΕ4 --- Εμπειρική συμπεριφορά.] Πώς συμπεριφέρονται επιλεγμένες
  παραγωγικές υλοποιήσεις σε ελεγχόμενο περιβάλλον μικρής συστάδας ως
  προς καθυστέρηση, \en{throughput} και ανοχή σε σφάλματα; Η εργασία
  δίνει δικές της μετρήσεις γι' αυτό στο κεφάλαιο της αποτίμησης.
  \item[ΕΕ5 --- Κλιμάκωση.] Πώς μεταβάλλεται η συμπεριφορά καθώς μεγαλώνει
  το \en{quorum}; Η εργασία προσεγγίζει το ερώτημα σε δύο επίπεδα, και τα
  δηλώνει χωριστά. Το πειραματικό επίπεδο καλύπτει μόνο μικρά μεγέθη
  συστάδας, όπου ελέγχεται το κόστος της αριθμητικής του
  Σχήματος~\ref{fig:quorum}. Το βιβλιογραφικό επίπεδο καλύπτει τα μεγάλα
  μεγέθη, όπου εμφανίζεται η ασυμπτωτική συμπεριφορά. Η εργασία δίνει
  δικές της μετρήσεις για το πειραματικό επίπεδο στο κεφάλαιο της
  αποτίμησης.
\end{description}

%%blk intro.aims.02 from=dfce2af38acb
Η εργασία δηλώνει τα όρια εδώ και όχι στα συμπεράσματα. Πρώτον, η εργασία
αποτιμά \emph{υλοποιήσεις} και όχι πρωτόκολλα. Το μετρήσιμο αντικείμενο
είναι πάντοτε ένα σύστημα με συγκεκριμένη έκδοση και συγκεκριμένη
ρύθμιση. Κάθε ισχυρισμός για το υποκείμενο πρωτόκολλο περνά υποχρεωτικά
από το ΕΕ3.

%%blk intro.aims.03 from=d606b8a15bc8
Δεύτερον, ο διαθέσιμος προϋπολογισμός υποδομής περιορίζει το εύρος των
μεγεθών συστάδας που μπορεί να μελετήσει ένα πείραμα. Το εύρος αυτό είναι
μικρό σε σχέση με τις κλίμακες όπου εμφανίζονται οι ενδιαφέρουσες
ασυμπτωτικές συμπεριφορές. Αυτός είναι ο λόγος του διαχωρισμού του ΕΕ5 σε
δύο επίπεδα.

%%blk intro.aims.04 from=67089110cff9
Τρίτον, ένα πρωτόκολλο μπορεί να καλυφθεί \emph{αναλυτικά} χωρίς δικές
μας μετρήσεις. Όπου συμβαίνει αυτό, το κεφάλαιο το δηλώνει στην αρχή
του. Οι όποιοι αριθμοί προέρχονται εκεί από τρίτες δημοσιεύσεις που
κατονομάζονται ρητά, και δεν στέκουν ποτέ στον ίδιο πίνακα με δικές μας
μετρήσεις. Το πλήρες πλαίσιο εγκυρότητας, τα μεθοδολογικά κριτήρια και οι
απειλές εγκυρότητας ανήκουν στο κεφάλαιο της μεθοδολογίας. Το παρόν
κεφάλαιο δεν τα προεξοφλεί.

%%blk intro.structure.head from=330c145eb75d
\section{Συνεισφορά και δομή της εργασίας}
\label{sec:structure}

%%blk intro.structure.01 from=241b572f4703
Η εργασία επιδιώκει τέσσερις συνεισφορές.

%%blk intro.structure.list from=4b7e7f203919
\begin{enumerate}[label=\textbf{(\alph*)},leftmargin=2.4em]
  \item Μια χρονολογημένη αφήγηση της εξέλιξης των πρωτοκόλλων, με ρητό
  κανόνα ταξινόμησης. Η αφήγηση αποδίδει κάθε μετάβαση σε κατονομασμένη
  ανταλλαγή.
  \item Μια ενιαία δομική περιγραφή. Το ίδιο σχήμα ανάγνωσης καλύπτει τα
  πρωτόκολλα και των δύο μοντέλων σφαλμάτων.
  \item Μια συγκριτική αποτίμηση επιλεγμένων παραγωγικών υλοποιήσεων σε
  ελεγχόμενο περιβάλλον, με προβλέψεις δηλωμένες εκ των προτέρων
  (\en{pre-registered predictions}). Η εργασία δηλώνει ρητά και τι
  \emph{δεν} μπορεί να συναχθεί από τα δεδομένα. Η εργασία δίνει δικές
  της μετρήσεις γι' αυτό στο κεφάλαιο της αποτίμησης.
  \item Ένα πλαίσιο απόφασης που αντιστοιχίζει απαιτήσεις σε οικογένειες
  πρωτοκόλλων.
\end{enumerate}

%%blk intro.structure.02 from=ec1f8fc7af8b
Η διάρθρωση ακολουθεί δύο χρονολογικές διαδρομές, μία ανά μοντέλο
σφαλμάτων, αντί για μία αυστηρά ενιαία. Η επιλογή είναι συνειδητή
απόκλιση από την αυστηρά χρονολογική παρουσίαση και τίθεται ρητά υπόψη
της επιβλέπουσας. Με ενιαία σειρά, η ακολουθία του
Σχήματος~\ref{fig:timeline} υποχρεώνει τον αναγνώστη να εναλλάξει μοντέλο
σφαλμάτων πέντε φορές. Το αποτέλεσμα είναι να παρατίθενται δίπλα-δίπλα
μεγέθη που δεν συγκρίνονται.

%%blk intro.structure.03 from=1bfe64225865
Το Κεφάλαιο~\ref{ch:background} θεμελιώνει το θεωρητικό υπόβαθρο: μοντέλα
σφαλμάτων και συγχρονισμού, SMR, \en{quorums} και κριτήρια ορθότητας. Τα
κεφάλαια της πρώτης διαδρομής καλύπτουν την οικογένεια κατάρρευσης, και
τα κεφάλαια της δεύτερης τη βυζαντινή. Κάθε πρωτόκολλο ακολουθεί κοινό
πρότυπο πέντε υποενοτήτων:

%%blk intro.structure.template from=14f76f86d68c
\begin{enumerate}[label=\arabic*.,leftmargin=2.4em]
  \item ιστορικό πλαίσιο και πρωτογενής πηγή·
  \item ο αλγόριθμος·
  \item η υλοποίηση που αποτιμήθηκε, ή η αιτιολόγηση της απουσίας της·
  \item οι μετρήσεις·
  \item η ανάγνωση των αποτελεσμάτων έναντι της πρόβλεψης που είχε δηλωθεί
  εκ των προτέρων.
\end{enumerate}

%%blk intro.structure.04 from=162cf9484024
Ακολουθεί το κεφάλαιο της μεθοδολογίας και της εγκυρότητας. Έπειτα
έρχονται η συγκεντρωτική σύγκριση, το πλαίσιο απόφασης και, τέλος, τα
συμπεράσματα με τους περιορισμούς και τη μελλοντική εργασία.

%%blk intro.sources.head from=d131226a5c4a
\section{Σημείωση για την ποιότητα των πηγών}
\label{sec:sources}

%%blk intro.sources.01 from=db5ff35becc9
Μεγάλο μέρος της βιβλιογραφίας των πρωτοκόλλων συναίνεσης κυκλοφορεί ως
\en{preprint}, ως \en{whitepaper} ή ως τεχνική τεκμηρίωση έργου. Γι' αυτό
τον λόγο η εργασία εφαρμόζει τρεις κανόνες σε όλα της τα κεφάλαια.
Πρώτον, κάθε αλγόριθμος περιγράφεται από την πρωτογενή του πηγή. Οι
ανασκοπήσεις και οι ενιαίες πλατφόρμες χρησιμοποιούνται για ενοποιητικά
σχήματα και για ταξινομήσεις. Όπου μια τέτοια εργασία μεταφέρει φράγμα ή
αριθμό τρίτου, πράγμα που συμβαίνει μόνο κατ' εξαίρεση, κατονομάζεται και
η πρωτογενής πηγή, όπως στη λεζάντα του Σχήματος~\ref{fig:quorum}.

%%blk intro.sources.02 from=d0da0f8f2fc7
Δεύτερον, οι μη κριθείσες πηγές δηλώνονται ως τέτοιες στο σημείο όπου
χρησιμοποιούνται. Η σήμανση \grey\ δείχνει τη γκρίζα βιβλιογραφία:
\en{whitepapers}, KIP και ιστοσελίδες. Η σήμανση \prep\ δείχνει τα
\en{preprints} και τις εκτεταμένες περιλήψεις χωρίς κρίση. Μια τρίτη
σήμανση, \unver, λέει κάτι διαφορετικό από τις δύο προηγούμενες: λέει ότι
τα βιβλιογραφικά στοιχεία της ίδιας της πηγής δεν ήταν δυνατόν να
ελεγχθούν ούτε στο άρθρο ούτε σε ευρετήριο, και πάρθηκαν από τον κατάλογο
αναφορών άλλης εργασίας. Ο αναγνώστης που θέλει να ακολουθήσει την
παραπομπή πρέπει να το ξέρει πριν το επιχειρήσει. Τρίτον, κάθε
δημοσιευμένος αριθμός που αναπαράγει η εργασία συνοδεύεται από το
περιβάλλον στο οποίο παρήχθη: υλικό, μοντέλο φόρτου, τοπολογία ανάπτυξης,
μέγεθος συστάδας και πλήθος επαναλήψεων. Οι λεζάντες των
Σχημάτων~\ref{fig:quorum} και~\ref{fig:scaling} το κάνουν αυτό. Αριθμός
χωρίς περιβάλλον δεν είναι τεκμήριο.

%%blk intro.sources.03 from=1662413bd872
Η ίδια αρχή εφαρμόζεται και στα βιβλιογραφικά στοιχεία. Στο αρχείο
βιβλιογραφίας κάθε εγγραφή φέρει σήμανση για το πώς επαληθεύτηκε: από την
πρώτη σελίδα του ίδιου του άρθρου, από αυθεντικό ευρετήριο μέσω του DOI,
ή καθόλου. Στην τελευταία περίπτωση η εγγραφή δηλώνεται ως μη
επαληθευμένη. Η διάκριση δεν είναι τυπολατρική. Τρεις περιπτώσεις όπου η
σήμανση απέτρεψε λάθος αξίζει να καταγραφούν.

%%blk intro.sources.04 from=f76fb1b59cec
Η πρώτη περίπτωση είναι το \cite{berger2023simulation}. Το αντίτυπο που
συλλέχθηκε είναι \en{preprint} και θα είχε δηλωθεί ως τέτοιο, ενώ το έργο
έχει κριθείσα εκδοχή σε συνέδριο. Η σήμανση \emph{επαλήθευσης} περιγράφει
το αντίτυπο που κρατάμε. Οι σημάνσεις \prep\ και \grey\ περιγράφουν το
\emph{έργο}. Εδώ τα δύο αποκλίνουν, και η απόκλιση εντοπίζεται μόνο με
έλεγχο ευρετηρίου. Το αντίτυπο είναι \en{preprint}, το έργο όχι, άρα το
\cite{berger2023simulation} \emph{δεν} φέρει \prep.

%%blk intro.sources.05 from=3e325a7e9e8a
Η δεύτερη περίπτωση είναι το \cite{oki1988viewstamped}. Καταχωρείται στο
Crossref με διαφορετική διατύπωση τίτλου από εκείνη των πρακτικών, και η
πλήρης μορφή επιβεβαιώνεται από τη βιβλιογραφία του ίδιου του
\cite{lamport1998parttime}. Η τρίτη περίπτωση είναι το
\cite{jalalzai2023fasthotstuff}. Ο τίτλος του είναι «\en{\dots A Fast and
Robust BFT Protocol}», όχι «\en{Resilient}», όπως εμφανίζεται σε
ορισμένους καταλόγους.

%%blk intro.sources.06 from=8fdfc7517d46
Μία τέταρτη περίπτωση δείχνει το ίδιο σε επίπεδο χαρακτήρα. Το
\en{corpus} έχει \emph{δύο} αντίτυπα του \cite{lamport1998parttime}, την
έκδοση της ACM και τη διορθωμένη του συγγραφέα. Τυπώνουν την ίδια
πληροφορία με άλλα γράμματα: «\en{Received: January 1990; accepted: March
1998}» η πρώτη, «\en{Received January 1990; Accepted March 1998}» η
δεύτερη. Όποιο και να παρατεθεί ως αυτούσιο, θα είναι λανθασμένο για το
άλλο. Γι' αυτό η λεζάντα του Σχήματος~\ref{fig:timeline} δίνει το
\emph{περιεχόμενο}. Όταν δύο αντίτυπα διαφέρουν, το αυτούσιο παράθεμα
απαιτεί να δηλωθεί και το αντίτυπο.

% =====================================================================
%  ΚΕΦΑΛΑΙΟ 2 — ΘΕΩΡΗΤΙΚΟ ΥΠΟΒΑΘΡΟ
%  Κανόνες που ισχύουν εδώ (Κεφ. 1 \S1.6):
%   - πρωτογενής πηγή για κάθε φράγμα και κάθε ορισμό
%   - καμία αυτούσια παράθεση από άρθρο που δεν βρίσκεται στο corpus/
%   - το πλαίσιο εγκυρότητας ΔΕΝ μπαίνει εδώ. Ανήκει στο κεφάλαιο της
%     μεθοδολογίας, όπως υπόσχεται το \S1.4.
%
%  Η ΑΓΓΛΙΚΗ ΕΙΝΑΙ ΤΟ ΠΡΩΤΟΤΥΠΟ. Αυτό το αρχείο ακολουθεί το
%  content/en/ch2.tex, μπλοκ προς μπλοκ.
% =====================================================================

%%blk background.head from=d5c676003f4b
\chapter{Θεωρητικό υπόβαθρο: μοντέλα, φράγματα και κριτήρια ορθότητας}
\label{ch:background}

%%blk background.intro.01 from=23b735575924
Το Κεφάλαιο~\ref{ch:intro} έδειξε ότι κάθε πρωτόκολλο συναίνεσης είναι
μια θέση σε έναν χώρο σχεδίασης. Αυτό το κεφάλαιο ορίζει τους άξονες του
χώρου. Χωρίς αυτούς, οι συγκρίσεις των επόμενων κεφαλαίων δεν έχουν
μονάδες.

%%blk background.intro.02 from=8545cc0077b9
Η σειρά είναι σκόπιμη. Πρώτα έρχεται το περιβάλλον στο οποίο τρέχει ένα
πρωτόκολλο, δηλαδή ο συγχρονισμός και τα σφάλματα. Έπειτα έρχεται το
πρόβλημα που λύνει. Έπειτα έρχονται τα φράγματα που κανένα πρωτόκολλο
δεν παραβιάζει. Τέλος έρχονται τα κριτήρια με τα οποία κρίνεται η
ορθότητα και η επίδοσή του.

%%blk background.synchrony.head from=751c57e4c1a0
\section{Δύο άξονες συγχρονισμού, όχι ένας}
\label{sec:synchrony}

%%blk background.synchrony.01 from=431e54fae8e8
Η βιβλιογραφία συχνά μιλά για «σύγχρονα» και «ασύγχρονα» συστήματα σαν
να υπάρχει ένας άξονας. Οι Dwork, Lynch και Stockmeyer χωρίζουν δύο
\cite{dwork1988consensus}:

%%blk background.synchrony.axes from=0339863ff5da
\begin{enumerate}[label=\textbf{(\alph*)},leftmargin=2.4em]
  \item \textbf{Η επικοινωνία.} Υπάρχει άνω φράγμα $\Delta$ στον χρόνο
  παράδοσης ενός μηνύματος;
  \item \textbf{Οι επεξεργαστές.} Υπάρχει άνω φράγμα $\Phi$ στον λόγο με
  τον οποίο το ρολόι ενός κόμβου τρέχει γρηγορότερα από ενός άλλου;
\end{enumerate}

%%blk background.synchrony.02 from=060fae9782ef
Ένα σύστημα είναι \emph{σύγχρονο} όταν και τα δύο φράγματα υπάρχουν και
είναι γνωστά εκ των προτέρων. Είναι \emph{ασύγχρονο} όταν \emph{δεν
υπάρχει κανένα} από τα δύο. Οι δύο άξονες είναι ανεξάρτητοι, άρα
ανάμεσά τους υπάρχουν και μεικτά μοντέλα, όπου η μία διάσταση είναι
σύγχρονη και η άλλη όχι. Ο διαχωρισμός δεν είναι σχολαστικισμός.
Ακριβώς γι' αυτόν τον λόγο οι ίδιοι συγγραφείς δίνουν πίνακα
ανθεκτικότητας με πέντε στήλες αντί για δύο, και τρεις από αυτές
αναπαράγει ο Πίνακας~\ref{tab:dls}.

%%blk background.synchrony.sub.partial from=ac00c81ae549
\subsection{Οι δύο μορφές μερικής συγχρονίας}

%%blk background.synchrony.03 from=60a8c8ba8645
Ανάμεσα στα δύο άκρα βρίσκεται η \emph{μερική συγχρονία}. Το ίδιο άρθρο
την ορίζει με δύο τρόπους, και η διαφορά τους έχει σημασία για την
εργασία.

%%blk background.synchrony.04 from=1ab78664fa06
Στην \textbf{πρώτη μορφή} το φράγμα $\Delta$ υπάρχει, αλλά δεν το
ξέρουμε. Ζητάμε πρωτόκολλο που λειτουργεί σωστά όποια κι αν είναι η
τιμή, δηλαδή πρωτόκολλο που δεν έχει το $\Delta$ ενσωματωμένο.

%%blk background.synchrony.05 from=4da263bdf928
Στη \textbf{δεύτερη μορφή} ξέρουμε το $\Delta$, αλλά το σύστημα
μηνυμάτων δεν το σέβεται πάντοτε. Οι συγγραφείς προσθέτουν έναν
περιορισμό. Σε κάθε εκτέλεση υπάρχει μια χρονική στιγμή που οι κόμβοι
δεν γνωρίζουν. Από εκείνη τη στιγμή και μετά, το σύστημα μηνυμάτων
σέβεται το $\Delta$. Την ονομάζουν \en{global stabilization time} και τη
συντομογραφούν GST.

%%blk background.synchrony.06 from=ac2470f8b4e9
Οι ίδιες δύο μορφές ορίζονται και για το $\Phi$. Γι' αυτόν τον λόγο ο
πίνακας του άρθρου έχει χωριστή στήλη για μερικώς σύγχρονους
επεξεργαστές με σύγχρονη επικοινωνία. Η εργασία χρησιμοποιεί τη δεύτερη
μορφή και για τους δύο άξονες, όπως και τα πρωτόκολλα που αποτιμά.

%%blk background.synchrony.fig.synchrony from=dc51e1fb5fbf
\dgfigure{fig:synchrony}%
{Τα τρία μοντέλα συγχρονισμού, με μία κάθετη γραμμή ανά μήνυμα.}%
{Τα τρία μοντέλα συγχρονισμού, με μία κάθετη γραμμή ανά μήνυμα. Στο
σύγχρονο μοντέλο κάθε μήνυμα φτάνει εντός γνωστού $\Delta$. Στο
ασύγχρονο δεν υπάρχει φράγμα. Στο μερικώς σύγχρονο υπάρχει άγνωστη
στιγμή GST, μετά την οποία ισχύει το $\Delta$ \cite{dwork1988consensus}.
Η GST είναι \textbf{μοντέλο ανάλυσης} και όχι μηχανισμός. Κανένας
κώδικας δεν την υπολογίζει. Ο όρος ανήκει στο ίδιο το άρθρο του 1988 και
δεν είναι μεταγενέστερος.}%
{% fig-synchrony — the three timing models, one vertical stroke per message.
% Structure figure: dgflow. dgCFT marks a message that arrives inside a known
% bound, dgBFT one that does not — the standard's own key/bad pairing.
% The axis is 7.5 cm long so that the whole figure, row labels included, comes
% out at \dgwidth like every other figure; the message positions are the
% original ones scaled by 7.5/9.6.
\begin{dgflow}
  % --- synchronous -------------------------------------------------------
  \draw[dgaxis] (0,3.0) -- (7.5,3.0);
  \node[anchor=east, inner sep=4pt] at (-0.1,3.0) {\fl{sync.model.synchronous}};
  \node[dgnote, anchor=north west] at (0,2.84) {\fl{sync.note.synchronous}};
  \foreach \x in {0.781,1.875,2.969,4.063,5.156,6.250}{
    \draw[dgarrow,-,draw=dgCFT] (\x,2.88) -- (\x,3.12);}

  % --- asynchronous ------------------------------------------------------
  \draw[dgaxis] (0,1.6) -- (7.5,1.6);
  \node[anchor=east, inner sep=4pt] at (-0.1,1.6) {\fl{sync.model.asynchronous}};
  \node[dgnote, anchor=north west] at (0,1.44) {\fl{sync.note.asynchronous}};
  \foreach \x in {0.469,0.859,2.344,5.078,5.313,6.719}{
    \draw[dgarrow,-,draw=dgBFT] (\x,1.48) -- (\x,1.72);}

  % --- partially synchronous --------------------------------------------
  \draw[dgaxis] (0,0.2) -- (7.5,0.2);
  \node[anchor=east, inner sep=4pt] at (-0.1,0.2) {\fl{sync.model.partial}};
  \foreach \x in {0.469,1.172,2.031,2.891}{
    \draw[dgarrow,-,draw=dgBFT] (\x,0.08) -- (\x,0.32);}
  \foreach \x in {4.063,4.922,5.781,6.641}{
    \draw[dgarrow,-,draw=dgCFT] (\x,0.08) -- (\x,0.32);}
  % The dashed line stops above the note: in the original it ran through it.
  \draw[dgaxis,-,dashed] (3.594,-0.45) -- (3.594,0.85);
  \node[dgtag, fill=dgPaper, text=dgAux, anchor=south] at (3.594,0.85) {\textbf{GST}};
  \node[dgtag, fill=dgPaper, text=dgBFT, anchor=north] at (1.62,-0.05)
       {\fl{sync.liveness.none}};
  \node[dgtag, fill=dgPaper, text=dgCFT, anchor=north] at (5.60,-0.05)
       {$\le\Delta$, \fl{sync.liveness.guaranteed}};
  \node[dgnote, anchor=north west] at (-1.2,-0.75) {\fl{sync.note.safety}};
\end{dgflow}
}

%%blk background.synchrony.07 from=31ff8b953ef1
Η δεύτερη μορφή είναι εκείνη που υιοθετούν τα βυζαντινά πρωτόκολλα της
εργασίας \cite{yin2019hotstuff} και \cite{malkhi2023hotstuff2}\prep. Η
πρακτική συνέπεια φαίνεται στο Σχήμα~\ref{fig:synchrony}. Η ασφάλεια
ισχύει σε ολόκληρο τον άξονα. Η ζωτικότητα εξασφαλίζεται μόνο δεξιά της
GST.

%%blk background.faults.head from=95886099d973
\section{Μοντέλα σφαλμάτων}
\label{sec:faults}

%%blk background.faults.01 from=b1275f514326
Ο δεύτερος άξονας ρωτά τι επιτρέπεται να κάνει ένας χαλασμένος κόμβος. Η
εργασία χρησιμοποιεί τέσσερις βαθμίδες, με αυξανόμενη κακοήθεια.

%%blk background.faults.02 from=5dc6adc7e0a5
Μια σύμβαση συμβόλων, μία φορά για όλο το κεφάλαιο. Τα $t$ και $f$
δηλώνουν το \emph{ίδιο} μέγεθος, το ανώτατο πλήθος ελαττωματικών κόμβων.
Το κεφάλαιο διατηρεί το σύμβολο κάθε πρωτογενούς πηγής, ώστε οι τιμές να
είναι αναγνωρίσιμες σε αυτήν.

%%blk background.faults.list from=ea920960c543
\begin{description}[leftmargin=2.2em,style=nextline,font=\normalfont\bfseries]
  \item[Κατάρρευση (\en{crash}).] Ο κόμβος σταματά και δεν ξαναμιλά. Το
  σφάλμα είναι \emph{απουσία} πληροφορίας. Το ειδικότερο μοντέλο
  \en{fail-stop} προσθέτει ότι οι υπόλοιποι κόμβοι αντιλαμβάνονται
  \emph{αξιόπιστα} τη διακοπή. Η διάκριση επανέρχεται στην
  Ενότητα~\ref{sec:smr}, όπου το φράγμα $t{+}1$ του Schneider
  προϋποθέτει \en{fail-stop} και όχι σκέτο \en{crash}.
  \item[Παραλείψεις (\en{omission}).] Ο κόμβος χάνει μηνύματα, αλλά όσα
  στέλνει είναι σωστά.
  \item[Βυζαντινό με αυθεντικοποίηση.] Ο κόμβος συμπεριφέρεται αυθαίρετα,
  αλλά δεν μπορεί να πλαστογραφήσει την υπογραφή άλλου κόμβου.
  \item[Βυζαντινό χωρίς αυθεντικοποίηση.] Καμία εγγύηση. Ο κόμβος λέει
  ό,τι θέλει σε όποιον θέλει, και συνεννοείται με τους υπόλοιπους
  ελαττωματικούς \cite{lamport1982byzantine}.
\end{description}

%%blk background.faults.03 from=716ec3fd47ac
Η συμπεριφορά που χωρίζει το βυζαντινό μοντέλο από τα υπόλοιπα έχει
όνομα. Είναι η \emph{αντιφατική δήλωση} (\en{equivocation}). Ο ίδιος
κόμβος λέει διαφορετικά πράγματα σε διαφορετικούς παραλήπτες, και οι
παραλήπτες δεν μπορούν να το διαπιστώσουν μόνοι τους. Το
Σχήμα~\ref{fig:equivocation} το δείχνει.

%%blk background.faults.fig.equivocation from=c119ec21a4c3
\dgfigure{fig:equivocation}%
{Κατάρρευση έναντι αντιφατικής δήλωσης: η μία συμπεριφορά που κάνει το
βυζαντινό μοντέλο ακριβό.}%
{Κατάρρευση έναντι αντιφατικής δήλωσης: η μία συμπεριφορά που κάνει το
βυζαντινό μοντέλο ακριβό. Αριστερά, ο $K_2$ σταματά. Οι $K_1$ και $K_3$
έχουν λιγότερη πληροφορία, αλλά όχι αντιφατική. Δεξιά, ο $K_2$ στέλνει
διαφορετική τιμή σε καθέναν και συνεχίζει να λειτουργεί. Ούτε ο $K_1$
ούτε ο $K_3$ μπορεί να το διαπιστώσει μόνος του. Το έμπιστο υλικό
μειώνει το φράγμα από $3f{+}1$ σε $2f{+}1$ ακριβώς επειδή εμποδίζει αυτή
τη συμπεριφορά. Η ταξινόμηση μεταφέρεται από την ανασκόπηση
\cite{amiri2024bedrock} (Παράρτ.~A, O7). Η πρωτογενής πηγή του φράγματος
είναι η γραμμή έργων έμπιστου υλικού στην οποία παραπέμπει η ανασκόπηση:
το TrInc, που το κατονομάζει, και άλλα στα οποία παραπέμπει με αριθμό
αναφοράς χωρίς να τα ονομάζει. Η εργασία \emph{δεν} τα έχει ελέγξει
ανεξάρτητα.}%
{% fig-equivocation — crash versus equivocation. Structure figure: dgflow.
% Left panel is the crash model, so every mark in it is dgCFT; right panel is
% the Byzantine model, so every mark in it is dgBFT. K1/K3 are ordinary nodes
% in both panels and stay dgbox grey. The two panels are 6 cm apart so that
% the pair comes out at \dgwidth.
\begin{dgflow}
  % ---- left: crash ------------------------------------------------------
  \node[text=dgCFT] at (1.5,3.35) {\textbf{\fl{equivocation.crash.title}}};
  \node[dgboxkey] (c2) at (0.0,1.2) {$K_2$};
  \node[dgbox]    (c1) at (3.0,2.4) {$K_1$};
  \node[dgbox]    (c3) at (3.0,0.0) {$K_3$};
  \draw[dgarrowkey] (c2) -- (c1)
        node[dgtag, fill=dgPaper, sloped, above, pos=0.5]
        {\fl{quote.open}$v=5$\fl{quote.close}};
  \draw[dgarrowkey] (c2) -- (c3)
        node[dgtag, fill=dgPaper, sloped, below, pos=0.5]
        {\fl{quote.open}$v=5$\fl{quote.close}};
  % K2 is gone: the cross sits just under the box it kills.
  \draw[dgarrowkey,-] (-0.30,0.65) -- (0.30,0.05);
  \draw[dgarrowkey,-] (-0.30,0.05) -- (0.30,0.65);
  \node[dgnote, anchor=north, align=center] at (1.5,-0.55)
       {\fl{equivocation.crash.note}};

  % ---- right: Byzantine -------------------------------------------------
  \node[text=dgBFT] at (7.5,3.35) {\textbf{\fl{equivocation.byzantine.title}}};
  \node[dgboxbad] (b2) at (6.0,1.2) {$K_2$};
  \node[dgbox]    (b1) at (9.0,2.4) {$K_1$};
  \node[dgbox]    (b3) at (9.0,0.0) {$K_3$};
  \draw[dgarrowbad] (b2) -- (b1)
        node[dgtag, fill=dgPaper, sloped, above, pos=0.5]
        {\fl{quote.open}$v=5$\fl{quote.close}};
  \draw[dgarrowbad] (b2) -- (b3)
        node[dgtag, fill=dgPaper, sloped, below, pos=0.5]
        {\fl{quote.open}$v=7$\fl{quote.close}};
  \node[dgnote, anchor=north, align=center] at (7.5,-0.55)
       {\fl{equivocation.byzantine.note}};
\end{dgflow}
}

%%blk background.faults.04 from=c7b1a25b883a
Ο Πίνακας~\ref{tab:dls} συνδυάζει τους δύο άξονες. Είναι το πιο
συμπυκνωμένο αποτέλεσμα του κεφαλαίου, και εξηγεί γιατί η μερική
συγχρονία δεν είναι δωρεάν.

%%blk background.faults.tab.dls from=42090924b75d
\begin{table}[t]
\centering
\caption[Το ελάχιστο πλήθος κόμβων ανά μοντέλο: το ελάχιστο $N$ για το
οποίο υπάρχει πρωτόκολλο συναίνεσης ανθεκτικό σε $t$ σφάλματα.]%
{Το ελάχιστο πλήθος κόμβων ανά μοντέλο: το ελάχιστο $N$ για το οποίο
υπάρχει πρωτόκολλο συναίνεσης ανθεκτικό σε $t$ σφάλματα. Ο πίνακας
αναπαράγει \textbf{τρεις από τις πέντε στήλες} του Πίνακα~I του
\cite{dwork1988consensus}. Η επιλογή αιτιολογείται στη σημείωση. Το
σύμβολο $\infty$ σημαίνει ότι δεν υπάρχει τέτοιο $N$, δηλαδή το FLP.
\emph{Δύο αναγνώσεις που αξίζουν:} στο σύγχρονο μοντέλο με
αυθεντικοποίηση ανέχεσαι \emph{οποιοδήποτε} πλήθος σφαλμάτων, και η
μερική συγχρονία ανεβάζει το κόστος του \en{fail-stop} από $t$ σε
$2t{+}1$, δηλαδή πληρώνεις για το ότι δεν ξέρεις πότε αρχίζει η
σταθερότητα.}
\label{tab:dls}
\small
\begin{tabular}{@{}l c c c@{}}
\toprule
\textbf{Είδος σφάλματος} & \textbf{Σύγχρονο} & \textbf{Ασύγχρονο} &
\textbf{Μερικώς σύγχρονο} \\
\midrule
\en{Fail-stop} (με την έννοια της πηγής) & $t$ & $\infty$ & $2t+1$ \\
Παραλείψεις               & $t$      & $\infty$ & $2t+1$ \\
Βυζαντινό με αυθεντικοποίηση & $t$   & $\infty$ & $3t+1$ \\
Βυζαντινό                 & $3t+1$   & $\infty$ & $3t+1$ \\
\bottomrule
\end{tabular}

\vspace{4pt}
\begin{minipage}{0.94\linewidth}\setstretch{1.08}\small
\textbf{Σημείωση ακρίβειας.} Η στήλη «Μερικώς σύγχρονο» αντιστοιχεί στην
περίπτωση όπου \emph{και} η επικοινωνία \emph{και} οι επεξεργαστές είναι
μερικώς σύγχρονοι. Το πρωτότυπο δίνει πέντε στήλες, επειδή διακρίνει και
τους συνδυασμούς όπου η μία διάσταση είναι σύγχρονη και η άλλη όχι. Η
έκδοση εδώ κρατά τις τρεις που χρησιμοποιεί η εργασία. Το ίδιο το άρθρο
βρίσκεται εκτός \en{corpus}. Το \emph{αντίτυπο αναφοράς} είναι το
ανοικτό θεσμικό αντίτυπο της ομάδας της Lynch στο MIT
(\texttt{groups.csail.mit.edu/tds/papers/Lynch/jacm88.pdf}). Από εκεί
διαβάσαμε τα μεγέθη του πίνακα, τους ορισμούς των $\Delta$ και $\Phi$,
τις δύο μορφές μερικής συγχρονίας και τον όρο GST. Δεν παρατίθεται
τίποτε αυτούσιο.

\smallskip
\textbf{Μία λέξη, δύο σημασίες.} Η πηγή ονομάζει την πρώτη γραμμή
\emph{\en{fail-stop}} και την ορίζει ως επεξεργαστή που εκτελεί σωστά,
μπορεί να σταματήσει οποτεδήποτε, και αφού σταματήσει δεν ξαναξεκινά.
\textbf{Δεν υπάρχει ρήτρα ανίχνευσης}. Δεν είναι αυτή η σημασία με την
οποία χρησιμοποιεί τη λέξη η Ενότητα~\ref{sec:faults} του κεφαλαίου,
ακολουθώντας το \cite{schneider1990implementing}, όπου το
\en{fail-stop} προσθέτει ότι οι υπόλοιποι κόμβοι αντιλαμβάνονται
\emph{αξιόπιστα} τη διακοπή. Η διαφορά κρίνει την ασύγχρονη στήλη: η
αξιόπιστη ανίχνευση είναι γνησίως ισχυρότερη υπόθεση από εκείνη που
κάνει ο πίνακας, και το $\infty$ της στήλης ανήκει στην ασθενέστερη
υπόθεση της πηγής. Η γραμμή επομένως φέρει το όνομα που της δίνει η
πηγή, και ο αναγνώστης πληροφορείται ποια από τις δύο σημασίες ισχύει.
Είναι η ίδια παγίδα με την Ενότητα~\ref{sec:smr}: μία λέξη, δύο
μετρήσεις.
\end{minipage}
\end{table}

%%blk background.smr.head from=30f775134ee8
\section{Συναίνεση, ατομική εκπομπή και SMR δεν είναι το ίδιο}
\label{sec:smr}

%%blk background.smr.01 from=3c65c3ca4e01
Τα τρία ονόματα χρησιμοποιούνται συχνά εναλλάξ. Δεν είναι συνώνυμα, και
η σύγχυση παράγει λάθος συγκρίσεις στα επόμενα κεφάλαια.

%%blk background.smr.list from=3299af9be4eb
\begin{description}[leftmargin=2.2em,style=nextline,font=\normalfont\bfseries]
  \item[Συναίνεση (\en{consensus}).] Οι κόμβοι συμφωνούν σε \emph{μία}
  τιμή.
  \item[Ατομική εκπομπή (\en{atomic broadcast}).] Οι κόμβοι παραδίδουν
  την \emph{ίδια ακολουθία} μηνυμάτων με την ίδια σειρά.
  \item[Αναπαραγωγή μηχανής καταστάσεων (SMR).] Η ατομική εκπομπή
  εφαρμόζεται σε ντετερμινιστική μηχανή, ώστε όλα τα αντίγραφα να
  καταλήξουν στην ίδια κατάσταση \cite{schneider1990implementing}.
\end{description}

%%blk background.smr.02 from=8216bb1b38e6
Το Σχήμα~\ref{fig:smr} συνοψίζει τη στοίβα. Η SMR είναι το επίπεδο στο
οποίο ζει η εργασία. Ο Schneider αναλύει την απαίτησή της σε δύο
συνθήκες. Κάθε μη ελαττωματικό αντίγραφο λαμβάνει κάθε αίτημα, και κάθε
μη ελαττωματικό αντίγραφο τα επεξεργάζεται με την ίδια σχετική σειρά. Η
πρώτη συνθήκη λέγεται \emph{συμφωνία}, η δεύτερη \emph{διάταξη}.
Πρόσεξε ότι η δεύτερη μιλά για \emph{σχετική} σειρά, και όχι για
καθολική χρονοσφραγίδα.

%%blk background.smr.03 from=4501fac53591
Πάνω από όλα αυτά στέκει μια υπόθεση που σπάει εύκολα στην πράξη. Η
ακολουθία των αιτημάτων πρέπει να καθορίζει πλήρως τις εξόδους της
μηχανής, και τίποτε άλλο δεν επιτρέπεται να τις επηρεάζει
\cite{schneider1990implementing}. Αν η μηχανή καλέσει γεννήτρια τυχαίων
αριθμών ή ρολόι, η SMR παύει να ισχύει.

%%blk background.smr.fig.smr from=8749db49f0a2
\dgfigure{fig:smr}%
{Η στοίβα: πού ζει το καθένα.}%
{Η στοίβα: πού ζει το καθένα. Η συναίνεση λύνει τη σειρά. Η μηχανή
καταστάσεων εκτελεί. Όλα τα πρωτόκολλα της εργασίας κατοικούν στο μεσαίο
κουτί, και γι' αυτό συγκρίνονται. Η ντετερμινιστικότητα των μηχανών
είναι υπόθεση και όχι συμπέρασμα \cite{schneider1990implementing}.}%
{% fig-smr — where each layer lives. Structure figure: dgflow.
% The middle box is the subject of the whole thesis, so it is dgboxkey; the
% clients and the replicas are ordinary components, so they are dgbox. Both
% CFT and BFT protocols live in the middle box, which is why it cannot carry
% either series colour.
\begin{dgflow}
  \node[dgbox]    (cl) at (0.0,0.0) {\fl{smr.clients}};
  \node[dgboxkey] (co) at (4.0,0.0)
       {\bfseries\fl{smr.consensus}\\\mdseries\fl{smr.agreement.order}};
  \node[dgbox]    (sm1) at (8.4, 0.95) {\fl{smr.machine} 1};
  \node[dgbox]    (sm2) at (8.4, 0.00) {\fl{smr.machine} 2};
  \node[dgbox]    (sm3) at (8.4,-0.95) {\fl{smr.machine} 3};
  \draw[dgarrow] (cl) -- (co);
  \draw[dgarrow] (co) -- (sm1);
  \draw[dgarrow] (co) -- (sm2);
  \draw[dgarrow] (co) -- (sm3);
  \node[dgnote, anchor=north, align=center] at (8.4,-1.45)
       {\fl{smr.deterministic}\\\fl{smr.same.io}};
  \node[dgnote, anchor=north] at (4.0,-1.15)
       {\fl{smr.here.live} Paxos, Raft, PBFT, HotStuff};
\end{dgflow}
}

%%blk background.smr.sub.trap from=d55660872921
\subsection{Μια παγίδα αρίθμησης: $2t{+}1$ και $3t{+}1$ δεν μετρούν το
ίδιο}

%%blk background.smr.04 from=7b142cd8d1bf
Ο Schneider γράφει ότι με βυζαντινά σφάλματα χρειάζονται τουλάχιστον
$2t{+}1$ αντίγραφα, ενώ με σφάλματα \en{fail-stop} αρκούν $t{+}1$
\cite{schneider1990implementing}. Η δεύτερη τιμή προϋποθέτει ότι η πτώση
είναι \emph{ανιχνεύσιμη}. Τότε αρκεί να πιστέψεις οποιονδήποτε ζωντανό
κόμβο. Το Κεφάλαιο~\ref{ch:intro} γράφει $3f{+}1$ για το βυζαντινό
μοντέλο. Τα δύο δεν συγκρούονται, επειδή μετρούν διαφορετικά πράγματα.

%%blk background.smr.trap.list from=51d196acccb0
\begin{itemize}
  \item Ο Schneider μετρά αντίγραφα των οποίων οι \emph{έξοδοι
  ψηφίζονται}. Με $2t{+}1$ αντίγραφα, η πλειοψηφία των εξόδων μένει
  σωστή ακόμη και μετά από $t$ σφάλματα. Η συμφωνία και η διάταξη
  θεωρούνται ήδη λυμένες.
  \item Το $3f{+}1$ είναι το φράγμα για να \emph{λυθεί} η συμφωνία,
  δηλαδή το φράγμα για το ίδιο το πρωτόκολλο συναίνεσης
  \cite{pease1980reaching}.
\end{itemize}

%%blk background.smr.05 from=ebe438b5427c
Η σύγχυση των δύο είναι εύκολη, επειδή και τα δύο μεγέθη γράφονται ως
συνάρτηση του πλήθους των σφαλμάτων. Η εργασία χρησιμοποιεί το $3f{+}1$
παντού, επειδή μετρά πρωτόκολλα συναίνεσης και όχι ψηφοφορία εξόδων.

%%blk background.bounds.head from=8abb95caf712
\section{Τι είναι αδύνατο και τι κοστίζει}
\label{sec:bounds}

%%blk background.bounds.roadmap from=dd5e7b13c920
Τέσσερα αποτελέσματα οριοθετούν κάθε πρωτόκολλο της εργασίας. Το πρώτο
λέει ότι σε πλήρως ασύγχρονο σύστημα κανένα ντετερμινιστικό πρωτόκολλο
δεν τερματίζει πάντοτε. Το δεύτερο δίνει το ελάχιστο πλήθος κόμβων που
ανέχεται δεδομένο πλήθος βυζαντινών σφαλμάτων. Το τρίτο δίνει το κόστος
σε μηνύματα. Το τέταρτο δίνει το ελάχιστο πλήθος βημάτων επικοινωνίας.
Παρουσιάζονται με αυτή τη σειρά, γιατί το καθένα στενεύει τον χώρο που
άφησε ανοιχτό το προηγούμενο.

%%blk background.bounds.sub.flp from=3e265ae78fed
\subsection{FLP}

%%blk background.bounds.01 from=734383c3efc2
Το θεμελιώδες αρνητικό αποτέλεσμα διατυπώνεται σε μία γραμμή. Οι
Fischer, Lynch και Paterson αποδεικνύουν ότι «\en{No consensus protocol
is totally correct in spite of one fault}» \cite{fischer1985impossibility}.
Το πρωτότυπο βρίσκεται στο \en{corpus} και η φράση παρατίθεται αυτούσια.

%%blk background.bounds.02 from=ed658280bad5
Η μία γραμμή όμως κουβαλά δύο υποθέσεις, και το ίδιο το άρθρο τις
δηλώνει ρητά. Το σύστημα είναι \emph{πλήρως ασύγχρονο} και το πρωτόκολλο
\emph{ντετερμινιστικό}. Η δεύτερη υπόθεση είναι η πόρτα εξόδου. Τα
τυχαιοποιημένα πρωτόκολλα δεν εμπίπτουν στο θεώρημα, και το ίδιο το
άρθρο παραπέμπει σε τέτοια εργασία. Όλα τα πρωτόκολλα της εργασίας είναι
ντετερμινιστικά, άρα το θεώρημα τα δεσμεύει.

%%blk background.bounds.03 from=36f70b480c6b
Τρεις ακρίβειες αλλάζουν τη διατύπωση του αποτελέσματος. Και τις τρεις
τις διαβάσαμε στο ίδιο το άρθρο.

%%blk background.bounds.flp.list from=a7b9801c9c72
\begin{enumerate}
  \item Το άρθρο \textbf{δεν} ορίζει την ορθότητα μέσω των επώνυμων
  ιδιοτήτων «συμφωνία» και «εγκυρότητα». Η λέξη \en{validity} δεν
  εμφανίζεται καθόλου, ενώ η \en{agreement} εμφανίζεται μόνο ως όνομα
  του προβλήματος. Ορίζει ένα πρωτόκολλο ως \emph{μερικώς ορθό} με δύο
  συνθήκες. Καμία προσβάσιμη διαμόρφωση δεν έχει πάνω από μία τιμή
  απόφασης, και για κάθε τιμή υπάρχει προσβάσιμη διαμόρφωση που την
  αποφασίζει. Η δεύτερη συνθήκη είναι η \emph{μη τετριμμένη}, που είναι
  ασθενέστερη από την εγκυρότητα.
  \item Το άρθρο \textbf{δεν} εξετάζει βυζαντινά σφάλματα, και το
  δηλώνει ρητά. Υποθέτει αξιόπιστο σύστημα μηνυμάτων και μία μόνο
  κατάρρευση. Αυτό κάνει το αποτέλεσμα \emph{ισχυρότερο}. Αν η συναίνεση
  είναι αδύνατη στο ασθενέστερο μοντέλο σφαλμάτων, είναι κατά μείζονα
  λόγο αδύνατη στο βυζαντινό.
  \item Η απαίτηση τερματισμού είναι \textbf{εξαιρετικά ασθενής}, και το
  άρθρο το δηλώνει. Αρκεί \emph{κάποιος} κόμβος να αποφασίσει τελικά,
  και όχι όλοι. Το αρνητικό αποτέλεσμα δεν οφείλεται λοιπόν ούτε σε
  κακόβουλους κόμβους ούτε σε χαμένα μηνύματα. Οφείλεται μόνο στο ότι η
  σιωπή δεν ξεχωρίζει από την αργοπορία.
\end{enumerate}

%%blk background.bounds.sub.third from=b7ca28913a8a
\subsection{Το φράγμα του ενός τρίτου}

%%blk background.bounds.04 from=0fc36e0f0a99
Το φράγμα $n \ge 3f+1$ για βυζαντινή συμφωνία χωρίς αυθεντικοποίηση
δημοσιεύεται πρώτη φορά το 1980, με τη μορφή του προβλήματος της
\emph{διαδραστικής συνέπειας} \cite{pease1980reaching}. Το όνομα λέει
την ίδια συνθήκη από την ανάποδη. Οι ελαττωματικοί πρέπει να είναι
λιγότεροι από το ένα τρίτο των κόμβων, άρα οι σωστοί ξεπερνούν τα δύο
τρίτα. Δύο χρόνια αργότερα οι ίδιοι συγγραφείς το ξαναδιατυπώνουν ως
αλληγορία των βυζαντινών στρατηγών \cite{lamport1982byzantine}.

%%blk background.bounds.05 from=81bbc88f29fa
Το ίδιο άρθρο του 1980 εξετάζει ήδη και την \emph{αυθεντικοποιημένη}
περίπτωση \cite{pease1980reaching}. Εκεί οι ελαττωματικοί κόμβοι μπορούν
να σιωπούν, αλλά δεν μπορούν να παραποιούν όσα αναμεταδίδουν. Το φράγμα
του ενός τρίτου δεν ισχύει σε αυτή την περίπτωση.

%%blk background.bounds.06 from=23bf0896781f
Η αλληγορία του 1982 προσθέτει την καθαρότερη διατύπωση. Ο αλγόριθμος
SM($m$) λύνει το πρόβλημα για οποιοδήποτε πλήθος στρατηγών, αρκεί η
υπογραφή ενός πιστού στρατηγού να μην πλαστογραφείται
\cite{lamport1982byzantine}. Το ίδιο άρθρο βάζει μια ρήτρα. Το πρόβλημα
είναι κενό με λιγότερους από $m{+}2$ στρατηγούς. Το αποτέλεσμα των
υπογραφών δεν είναι λοιπόν συνεισφορά του 1982, αλλά του 1980.

%%blk background.bounds.07 from=24627e844d97
Το φράγμα $3f{+}1$ που χρησιμοποιούν όλα τα πρωτόκολλα της εργασίας δεν
προέρχεται λοιπόν από την απουσία υπογραφών. Προέρχεται από τη
\emph{μερική συγχρονία}, όπως δείχνει η γραμμή «Βυζαντινό με
αυθεντικοποίηση» του Πίνακα~\ref{tab:dls}. Εκεί η απαίτηση είναι
$n \ge 3t{+}1$ ακόμη και όταν τα μηνύματα είναι υπογεγραμμένα.

%%blk background.bounds.quote from=d8bd4398be9d
\begin{quote}\small
Αν σε ρωτήσουν «γιατί οι υπογραφές δεν σε γλιτώνουν από το $3f{+}1$;», η
απάντηση είναι αυτή η γραμμή του πίνακα. Οι υπογραφές αρκούν στο
\emph{σύγχρονο} μοντέλο. Στο μερικώς σύγχρονο δεν αρκούν.
\end{quote}

%%blk background.bounds.sub.messages from=14ef45ab1fd7
\subsection{Το κόστος σε μηνύματα}

%%blk background.bounds.08 from=6a2ef96a77cd
Η αδυναμία δεν είναι το μόνο φράγμα. Υπάρχει και κάτω φράγμα στο πλήθος
των μηνυμάτων. Οι Dolev και Reischuk αποδεικνύουν ότι κάθε
\textbf{ντετερμινιστικός} αλγόριθμος βυζαντινής συμφωνίας που ανέχεται
$t$ σφάλματα έχει εκτέλεση στην οποία οι σωστοί κόμβοι στέλνουν
τετραγωνικό ως προς $t$ πλήθος μηνυμάτων \cite{dolev1985bounds}.

%%blk background.bounds.09 from=d5b89e6bdb84
Τρία σημεία απαιτούν προσοχή, και το κεφάλαιο δηλώνει και τα τρία εδώ
αντί να τα αποσιωπήσει.

%%blk background.bounds.messages.list from=503bf627cf56
\begin{itemize}
  \item Το θεώρημα \textbf{δεν} περιορίζεται στους μη
  αυθεντικοποιημένους αλγορίθμους. Η υπόθεσή του δεν περιέχει ρήτρα
  αυθεντικοποίησης, άρα καλύπτει και τις δύο περιπτώσεις. Πολλές
  δευτερογενείς πηγές το γράφουν λανθασμένα.
  \item Ο περιορισμός στον \textbf{ντετερμινισμό} δεν είναι
  διακοσμητικός. Τα τυχαιοποιημένα πρωτόκολλα μπορούν να πετύχουν
  υποτετραγωνική επικοινωνία, και πληρώνουν με πιθανότητα σφάλματος.
  Αυτό δεν ισχύει πάντα. Απέναντι σε \emph{ισχυρά προσαρμοστικό}
  αντίπαλο, που διαγράφει μηνύματα εκ των υστέρων, το τετραγωνικό φράγμα
  ισχύει και για αυτά \cite{abraham2019ccba}. Δεν σπάνε λοιπόν το
  φράγμα. Βγαίνουν από την υπόθεσή του, και μόνο υπό ασθενέστερο μοντέλο
  αντιπάλου. Όλα τα πρωτόκολλα της εργασίας είναι ντετερμινιστικά.
  \item Το ίδιο το άρθρο του περιοδικού βρίσκεται επί πληρωμή και εκτός
  \en{corpus}. Ελέγξαμε το περιεχόμενο στην προγενέστερη εκδοχή
  συνεδρίου, και η αρίθμηση των θεωρημάτων μπορεί να διαφέρει.
\end{itemize}

%%blk background.bounds.10 from=173331c5a7cb
Η συνέπεια αφορά άμεσα τα κεφάλαια των βυζαντινών πρωτοκόλλων. Όταν το
HotStuff διαφημίζει \emph{γραμμική} αλλαγή όψης, ο ισχυρισμός ισχύει και
για την αλλαγή όψης, εφόσον ο ηγέτης που αναλαμβάνει είναι σωστός και το
δίκτυο σταθερό \cite{yin2019hotstuff}. Το κάτω φράγμα δεν παραβιάζεται.
Μετατοπίζεται στη χειρότερη περίπτωση. Το ίδιο άρθρο αναφέρει $O(fn)$
\emph{αυθεντικοποιητές} πάνω σε $f$ διαδοχικούς ελαττωματικούς ηγέτες,
και αυτό είναι $O(n^{2})$ όταν το $f$ μεγαλώνει μαζί με το $n$.

%%blk background.bounds.11 from=4db334762ea8
\textbf{Προσοχή στη μονάδα.} Το HotStuff μετρά \en{authenticators}, ενώ
το φράγμα των Dolev και Reischuk μετρά \emph{μηνύματα}. Τα δύο μεγέθη
δεν ταυτίζονται, και το ίδιο το άρθρο του HotStuff το δηλώνει. Ένα
μήνυμα που κουβαλά $n$ υπογραφές μετρά όσο $n$ μηνύματα με μία υπογραφή
το καθένα. Η αντιπαράθεση εδώ είναι ποιοτική και όχι λογιστική.

%%blk background.bounds.sub.twostep from=4b21385912d2
\subsection{Το φράγμα των δύο βημάτων}

%%blk background.bounds.12 from=132bac2fe5d2
Η τελευταία κατηγορία φραγμάτων αφορά τον αριθμό των φάσεων. \emph{Δύο
βήματα} σημαίνει ότι μια εντολή δεσμεύεται με δύο μόνο διαδοχικές
φάσεις. Αυτές είναι η πρόταση του ηγέτη και ένας γύρος ψήφων.

%%blk background.bounds.13 from=8b82c5ddf6c7
\textbf{Προσοχή στη σύμβαση αρίθμησης.} Η ανασκόπηση
\cite{amiri2024bedrock} μετρά το PBFT με \emph{τρεις} φάσεις διάταξης
και το FaB με \emph{δύο}. Ο Πίνακας~\ref{tab:delays} ακολουθεί τη
σύμβαση του \cite{malkhi2023hotstuff2}\prep, όπου το PBFT εμφανίζεται με
δύο. Οι δύο αριθμήσεις δεν συγκρίνονται μεταξύ τους.

%%blk background.bounds.14 from=a21d7995b2e7
Η εξοικονόμηση της μίας φάσης κόστιζε αρχικά $5f{+}1$ κόμβους
\cite{martin2006fab}\unver. Με αυθεντικοποίηση αρκούν $5f{-}1$, και αυτό είναι
το κάτω φράγμα \cite{kuznetsov2021optimal,abraham2021goodcase}\unver.

%%blk background.bounds.15 from=a1c8a7aed969
Η ανασκόπηση \cite{amiri2024bedrock} μεταφέρει και τα δύο μεγέθη, και
από εκεί τα πήρε η εργασία. Το κείμενο κατονομάζει τις πρωτογενείς
πηγές, αλλά η εργασία δεν τις έχει ελέγξει ανεξάρτητα. Ο
Πίνακας~\ref{tab:bounds} το δηλώνει.

%%blk background.bounds.tab.bounds from=db1b365d08c0
\begin{table}[t]
\centering
\caption[Σύνοψη των φραγμάτων που χρησιμοποιεί η εργασία, με την
πρωτογενή τους πηγή.]%
{Σύνοψη των φραγμάτων που χρησιμοποιεί η εργασία, με την πρωτογενή τους
πηγή. Η στήλη «Τεκμηρίωση» είναι μέρος του τεκμηρίου και έχει τρεις
τιμές. \emph{Corpus}: το άρθρο υπάρχει και το διαβάσαμε.
\emph{Ευρετήριο}: τα βιβλιογραφικά επαληθεύτηκαν σε DOI ή σε ευρετήριο,
και το περιεχόμενο σε ανοικτό θεσμικό ή αυθεντικό αντίτυπο του
συγγραφέα, ή σε προγενέστερη εκδοχή. \emph{Μη επαληθευμένη}: τα στοιχεία
αντλήθηκαν από τον κατάλογο αναφορών του \cite{amiri2024bedrock}. Όπου η
στήλη δεν λέει \emph{corpus}, το κείμενο δεν παραθέτει τίποτε
αυτούσιο.}
\label{tab:bounds}
\small
\begin{tabular}{@{}l l c@{}}
\toprule
\textbf{Φράγμα} & \textbf{Πρωτογενής πηγή} & \textbf{Τεκμηρίωση} \\
\midrule
Αδυναμία σε ασύγχρονο, 1 κατάρρευση & \cite{fischer1985impossibility} & corpus \\
$n \ge 3f+1$, βυζαντινό & \cite{pease1980reaching} & ευρετήριο \\
Υπογραφές: το φράγμα καταρρέει (σύγχρονο) &
  \cite{pease1980reaching,lamport1982byzantine} & ευρετήριο \\
$n \ge 2f+1$ / $3f+1$ σε μερική συγχρονία & \cite{dwork1988consensus} & ευρετήριο \\
$\Omega(t^2)$ μηνύματα, ντετερμινιστικά & \cite{dolev1985bounds} & ευρετήριο \\
$5f{+}1$ για μία φάση λιγότερη & \cite{martin2006fab}\unver\ & μη επαληθευμένη \\
$5f{-}1$ κάτω φράγμα δύο βημάτων &
  \cite{kuznetsov2021optimal,abraham2021goodcase}\unver\ & μη επαληθευμένη \\
\bottomrule
\end{tabular}
\end{table}

%%blk background.bounds.16 from=0e496d95948a
Ο Πίνακας~\ref{tab:bounds} συγκεντρώνει τα φράγματα του κεφαλαίου μαζί
με την τεκμηρίωσή τους.

%%blk background.quorums.head from=8808e23da0f0
\section{Απαρτίες: η μία ιδιότητα που κάνει τη δουλειά}
\label{sec:quorums}

%%blk background.quorums.01 from=bb323c478274
Μια \emph{απαρτία} (\en{quorum}) είναι το ελάχιστο σύνολο κόμβων που
κάνει μια απόφαση οριστική. Η μοναδική ιδιότητα που χρειάζεται είναι η
τομή.

%%blk background.quorums.sub.derivation from=6a02fae17460
\subsection{Η παραγωγή}

%%blk background.quorums.02 from=f2d0985c53ca
Το Σχήμα~\ref{fig:quorum} του Κεφαλαίου~\ref{ch:intro} δίνει τα μεγέθη.
Εδώ είναι το επιχείρημα από το οποίο προκύπτουν, σε τρία βήματα.

%%blk background.quorums.list from=888b33260d40
\begin{enumerate}
  \item Δεν μπορείς να περιμένεις πάνω από $n-f$ απαντήσεις. Οι $f$
  ελαττωματικοί μπορεί να σιωπούν, και η σιωπή δεν ξεχωρίζει από την
  αργοπορία. Άρα η απαρτία έχει μέγεθος $q = n-f$.
  \item Δύο απαρτίες τέμνονται σε τουλάχιστον $2q - n = n - 2f$ κόμβους.
  \item Στο βυζαντινό μοντέλο η τομή πρέπει να περιέχει τουλάχιστον έναν
  \emph{τίμιο} κόμβο. Στην τομή κρύβονται έως $f$ βυζαντινοί, άρα
  $n - 2f > f$, δηλαδή $n \ge 3f+1$.
\end{enumerate}

%%blk background.quorums.03 from=4d420e6f4761
Στο μοντέλο κατάρρευσης το βήμα~3 δεν υπάρχει. Αρκεί $n - 2f > 0$,
δηλαδή $n \ge 2f+1$. Αυτό το ένα βήμα είναι ολόκληρη η διαφορά κόστους
μεταξύ των δύο μοντέλων.

%%blk background.quorums.quote from=6436242bcf09
\begin{quote}\small
\textbf{Τι είναι και τι δεν είναι αυτό.} Η παραπάνω παραγωγή είναι
επιχείρημα \emph{σχεδίασης απαρτίας}. Δείχνει ότι οι απαρτίες μεγέθους
$n-f$ χρειάζονται $n \ge 3f+1$. Δεν είναι η γενική απόδειξη
αδυνατότητας, που ανήκει στο \cite{pease1980reaching}.
\end{quote}

%%blk background.quorums.sub.intersection from=34a4b93980f6
\subsection{Η τομή απαιτείται λιγότερο απ' όσο νομίζαμε}

%%blk background.quorums.04 from=e1b30045befd
Από το Paxos και έπειτα, η πλειοψηφική απαρτία θεωρούνταν αναγκαία επί
Οι Howard, Malkhi και Spiegelman δείχνουν ότι
δεν είναι \cite{howard2016flexiblepaxos}. Στο Paxos, η τομή απαιτείται
μόνο \emph{μεταξύ} των δύο φάσεων. Οι απαρτίες της Φάσης~1 δεν
χρειάζεται να τέμνονται μεταξύ τους, ούτε εκείνες της Φάσης~2 μεταξύ
τους.

%%blk background.quorums.05 from=38f7d807d120
Η πρακτική συνέπεια είναι ότι μπορείς να μικρύνεις την απαρτία της
Φάσης~2, δηλαδή της σταθερής κατάστασης. Πληρώνεις με μεγαλύτερη απαρτία
στην εκλογή ηγέτη. Το κεφάλαιο του Paxos επανέρχεται σε αυτό, επειδή
αλλάζει την ανάγνωση της αριθμητικής του Σχήματος~\ref{fig:quorum}.

%%blk background.correctness.head from=7b9e028490ea
\section{Κριτήρια ορθότητας}
\label{sec:correctness}

%%blk background.correctness.roadmap from=1f66b7a0fe6c
Δύο ερωτήματα κρίνουν αν μια αναπαραγόμενη υπηρεσία είναι ορθή, και
είναι ξεχωριστά ερωτήματα. Το πρώτο είναι τι εγγυάται το πρωτόκολλο για
τη σειρά που χτίζει· αυτό είναι η ασφάλεια και η ζωντάνια. Το δεύτερο
είναι τι δικαιούται να παρατηρήσει ο πελάτης από έξω· αυτό είναι η
γραμμικοποιησιμότητα. Η ενότητα τα παίρνει με αυτή τη σειρά, γιατί το
δεύτερο διατυπώνεται με όρους του πρώτου.

%%blk background.correctness.sub.safety from=1273024c3aa1
\subsection{Ασφάλεια και ζωτικότητα}

%%blk background.correctness.01 from=fa7fec866fba
Ένα πρωτόκολλο κρίνεται από δύο κατηγορίες ιδιοτήτων. Η \emph{ασφάλεια}
λέει ότι κάτι κακό δεν συμβαίνει ποτέ. Η \emph{ζωτικότητα} λέει ότι κάτι
καλό συμβαίνει τελικά.

%%blk background.correctness.02 from=a8a76350284b
Η ασυμμετρία που διέπει όλα τα πρωτόκολλα της εργασίας είναι απλή. Η
ασφάλεια διατηρείται πάντοτε, ακόμη και σε πλήρη διαμέριση δικτύου. Η
ζωτικότητα θυσιάζεται. Η επιλογή είναι σχεδιαστική και όχι ατύχημα. Ένα
σύστημα που σταματά διορθώνεται, ενώ ένα σύστημα που έγραψε δύο
διαφορετικές τιμές δεν διορθώνεται.

%%blk background.correctness.sub.lin from=f31ac030bd03
\subsection{Γραμμικοποιησιμότητα}

%%blk background.correctness.03 from=4d58102b4893
Οι Herlihy και Wing ορίζουν το κριτήριο που περιμένει ο πελάτης
\cite{herlihy1990linearizability}. Κάθε πράξη φαίνεται να συμβαίνει
ακαριαία σε κάποιο σημείο ανάμεσα στην κλήση της και στην απόκρισή της.
Οι πράξεις που δεν επικαλύπτονται χρονικά διατηρούν τη σειρά τους.

%%blk background.correctness.04 from=6ee23d52022b
Το δεύτερο σκέλος κάνει όλη τη δουλειά. Είναι περιορισμός
\emph{πραγματικού χρόνου}, και γι' αυτό η γραμμικοποιησιμότητα κοστίζει
συντονισμό.

%%blk background.correctness.05 from=d395b473ce06
Δύο ιδιότητες τη χωρίζουν από τη σειριοποιησιμότητα των βάσεων
δεδομένων, και οι δύο έχουν σημασία για την εργασία.

%%blk background.correctness.list from=7d5c69fe512d
\begin{description}[leftmargin=2.2em,style=nextline,font=\normalfont\bfseries]
  \item[Τοπικότητα (\en{locality}).] Ένα σύστημα είναι γραμμικοποιήσιμο
  αν κάθε αντικείμενό του είναι. Η ιδιότητα \emph{συντίθεται}. Ούτε η
  σειριοποιησιμότητα ούτε η αυστηρή σειριοποιησιμότητα είναι τοπικές
  \cite{herlihy1990linearizability}.
  \item[Μη φραγή (\en{nonblocking}).] Μια εκκρεμής κλήση ολικής πράξης
  δεν υποχρεώνεται ποτέ να περιμένει άλλη εκκρεμή κλήση. Η
  σειριοποιησιμότητα είναι εγγενώς φράσσουσα
  \cite{herlihy1990linearizability}.
\end{description}

%%blk background.correctness.06 from=5c3e44100634
Οι ίδιοι συγγραφείς περιγράφουν τη γραμμικοποιησιμότητα ως ειδική
περίπτωση της αυστηρής σειριοποιησιμότητας, όπου κάθε δοσοληψία
περιορίζεται σε μία πράξη πάνω σε ένα αντικείμενο.

%%blk background.performance.head from=3c4350eec769
\section{Τι σημαίνει «επίδοση» για ένα πρωτόκολλο}
\label{sec:performance}

%%blk background.performance.01 from=7df6d5147d38
Αυτή η ενότητα είναι σύντομη και φέρει βάρος. Ορίζει, μία φορά, τη
διαφορά ανάμεσα σε δύο μεγέθη που το Κεφάλαιο~\ref{ch:intro} έδειξε ότι
συγχέονται.

%%blk background.performance.list from=c4e4829ca152
\begin{description}[leftmargin=2.2em,style=nextline,font=\normalfont\bfseries]
  \item[Φάσεις.] Πόσοι διαδοχικοί γύροι πρότασης και ψήφου χρειάζονται
  για να δεσμευτεί μια εντολή. Είναι ιδιότητα του \emph{πρωτοκόλλου} και
  \emph{συγκρίσιμη} μεταξύ άρθρων, εφόσον δηλωθεί η σύμβαση αρίθμησης.
  \item[Καθυστερήσεις μηνυμάτων.] Πόσες διαδοχικές αποστολές
  χρειάζονται για να δεσμευτεί μια εντολή. Στο HotStuff μία φάση
  αντιστοιχεί σε δύο καθυστερήσεις. Αλλού η αντιστοιχία διαφέρει, και
  γι' αυτό δηλώνεται.
  \item[Πραγματικός χρόνος.] Πόσα χιλιοστά του δευτερολέπτου
  μεσολαβούν. Εξαρτάται από την υλοποίηση, το υλικό και την τοπολογία
  του δικτύου. Δεν είναι συγκρίσιμος μεταξύ άρθρων. Η εργασία δίνει
  δικές της μετρήσεις πραγματικού χρόνου στο κεφάλαιο της αποτίμησης.
\end{description}

%%blk background.performance.02 from=77bc94736876
Ο Πίνακας~\ref{tab:delays} δίνει το πρώτο μέγεθος, σε \emph{φάσεις}, και
η σημείωσή του δηλώνει τη μετατροπή σε καθυστερήσεις. Είναι ο μόνος
πίνακας επίδοσης του κεφαλαίου, επειδή είναι ο μόνος που δεν εξαρτάται
από περιβάλλον.

%%blk background.performance.tab.delays from=7063ea00458b
\begin{table}[t]
\centering
\caption[Φάσεις στην ευτυχή διαδρομή: πόσες φάσεις χρειάζονται για μία
απόφαση, με σωστό ηγέτη και μετά την GST.]%
{Φάσεις στην ευτυχή διαδρομή: πόσες φάσεις χρειάζονται για μία απόφαση,
με σωστό ηγέτη και μετά την GST. Δεν υπάρχουν εδώ χιλιοστά του
δευτερολέπτου, και αυτό είναι το νόημα της ενότητας. Πηγές:
\cite{malkhi2023hotstuff2}\prep\ (Πίνακας~1) για τα τέσσερα βυζαντινά,
και \cite{ongaro2014raft} \S5.3 για τη γραμμή Multi-Paxos / Raft, όπου
σε σταθερή κατάσταση «\en{a new entry can be replicated with a single
round of RPCs to a majority of the cluster}».}
\label{tab:delays}
\small
\begin{tabular}{@{}l c c@{}}
\toprule
\textbf{Πρωτόκολλο} & \textbf{Φάσεις} & \textbf{Μοντέλο} \\
\midrule
Multi-Paxos / Raft & 1 & κατάρρευση \\
PBFT               & 2 & βυζαντινό \\
Tendermint         & 2 & βυζαντινό, μη αποκρίσιμο \\
HotStuff           & 3 & βυζαντινό \\
HotStuff-2\prep    & 2 & βυζαντινό \\
\bottomrule
\end{tabular}

\vspace{4pt}
\begin{minipage}{0.94\linewidth}\setstretch{1.08}\small
\textbf{Προσοχή στη λογιστική.} Ο πίνακας μετρά \emph{φάσεις}, με την
ίδια έννοια και την ίδια πηγή που χρησιμοποιεί ο
Πίνακας~\ref{tab:viewchange} του Κεφαλαίου~\ref{ch:intro}. Στο HotStuff
κάθε φάση αντιστοιχεί σε \emph{δύο} διαδοχικές καθυστερήσεις μηνυμάτων,
δηλαδή πρόταση του ηγέτη και ψήφοι. Αυτό το δηλώνει η ίδια η πηγή
\emph{για το HotStuff}, άρα οι τρεις φάσεις του είναι έξι καθυστερήσεις.
Η αντιστοιχία \textbf{δεν} είναι καθολική. Το PBFT εμφανίζεται με $2$
επειδή έχει δύο γύρους ψηφοφορίας μετά την πρόταση, δηλαδή τρεις
καθυστερήσεις συνολικά και όχι τέσσερις. Ο πίνακας συγκρίνει λοιπόν
\emph{φάσεις όπως τις μετρά η πηγή}, και όχι καθυστερήσεις. Η γραμμή
Multi-Paxos / Raft δεν προέρχεται από την ίδια πηγή και σημειώνεται
χωριστά στη λεζάντα. Ο αριθμός των φάσεων \textbf{δεν} μεταφράζεται σε
χρόνο χωρίς να δηλωθεί το δίκτυο. Τρεις φάσεις σε LAN κοστίζουν λιγότερο
από δύο σε δίκτυο ευρείας περιοχής.
\end{minipage}
\end{table}

%%blk background.performance.03 from=e25074d85178
Η ενότητα κλείνει με τον κανόνα που θα εφαρμοστεί σε κάθε επόμενο
κεφάλαιο. Όταν συγκρίνονται πρωτόκολλα, συγκρίνονται φάσεις. Όταν
συγκρίνονται υλοποιήσεις, συγκρίνονται χρόνοι --- και μόνο εφόσον
δηλωθεί το περιβάλλον. Το πλήρες πλαίσιο εγκυρότητας, τα επίπεδα
συγκρισιμότητας και οι απειλές εγκυρότητας αναπτύσσονται στο κεφάλαιο
της μεθοδολογίας, όπως υπόσχεται το \S1.4.


% --- Ενότητα 3η: ΤΕΛΙΚΑ ΣΤΟΙΧΕΙΑ -------------------------------------
% Δεν αριθμείται ως κεφάλαιο, αλλά μπαίνει στα περιεχόμενα (Οδηγίες,
% «Πίνακας Περιεχομένων»).
\cleardoublepage
\phantomsection
\addcontentsline{toc}{chapter}{Βιβλιογραφία}
\bibliographystyle{IEEEtran}
\bibliography{bib/references}

% Τα Παραρτήματα μπαίνουν ΜΕΤΑ τη Βιβλιογραφία: έτσι λένε οι Οδηγίες δύο
% φορές (Πίνακας Α.1 και «Παραρτήματα»). Το Πρότυπο Word τα βάζει πριν·
% υπερισχύει ο γραπτός κανόνας. Τιτλοφορούνται ΠΑΡΑΡΤΗΜΑ Α, Β, … και
% μπαίνουν στα περιεχόμενα. Για να προσθέσεις ένα: content/el/appA.tex
% με \chapter{Τίτλος}, βγάλε τα σχόλια από τις δύο γραμμές, και το ίδιο
% στο main-en.tex. ΠΡΩΤΑ διόρθωσε το .docx: το chapters() του
% docx/config.py παίρνει ΚΑΘΕ content/<lang>/*.tex ως κεφάλαιο, οπότε το
% appA.tex θα γινόταν Κεφάλαιο 1. Η τρίτη γραμμή δίνει ελληνικά γράμματα
% (Α, Β, Γ)· χωρίς αυτήν το numerals=arabic τυπώνει λατινικά A, B, C.
% \appendix
% \renewcommand{\thechapter}{\ifcase\value{chapter}\or Α\or Β\or Γ\or Δ\or Ε\fi}
% \include{content/el/appA}

\end{document}
